\documentclass[11pt,a4paper]{article}

\usepackage[a4paper,margin=2.3cm,headheight=15pt]{geometry}
\usepackage{fontspec}
\usepackage{amsmath,amssymb,amsthm,mathtools,bm}
\usepackage{booktabs,tabularx,longtable,array,multirow}
\usepackage{caption}
\usepackage{xcolor}
\usepackage{enumitem}
\usepackage{fancyhdr}
\usepackage{hyperref}
\usepackage{xurl}
\usepackage{listings}
\lstdefinestyle{spec}{
  language=C, basicstyle=\ttfamily\footnotesize, columns=fullflexible,
  keepspaces=true, breaklines=true, breakatwhitespace=true,
  showstringspaces=false, commentstyle=\color{black!55}\itshape,
  frame=single, framerule=0.4pt, rulecolor=\color{black!35},
  xleftmargin=0.5em, xrightmargin=0.5em, aboveskip=1.1em, belowskip=1.1em,
}
\lstset{style=spec}

\hypersetup{
  colorlinks=true,
  linkcolor=blue!55!black,
  citecolor=green!45!black,
  urlcolor=blue!65!black,
  pdftitle={A high-order meshless departure-point interpolation operator for
            semi-Lagrangian lattice Boltzmann methods},
  pdfauthor={Edit17 CCSLBM project}
}

\pagestyle{fancy}
\fancyhf{}
\lhead{\small LABFM departure-point interpolation}
\rhead{\small Theory and algorithm specification}
\cfoot{\thepage}

\setlength{\parindent}{1.2em}
\setlength{\parskip}{0.25em}
\setlength{\emergencystretch}{3em}
\setlist{nosep,leftmargin=2.0em}
\renewcommand{\arraystretch}{1.2}

% 可換行、左對齊的固定寬度欄（取代不換行的 l，避免把 X 欄擠扁）
\newcolumntype{P}[1]{>{\raggedright\arraybackslash}p{#1}}
% 寬表統一樣式：小一級字、放寬列距、允許連字斷行
\newenvironment{widetable}
  {\begingroup\small\setlength{\tabcolsep}{4pt}%
   \renewcommand{\arraystretch}{1.15}%
   \setlength{\LTpre}{0.8em}\setlength{\LTpost}{0.8em}}
  {\endgroup}

\newcommand{\rji}{\bm{r}_{ji}}
\newcommand{\New}{\textcolor{blue!60!black}{\textsf{[N]}}}
\newcommand{\TBD}{\textcolor{red!70!black}{\textsf{[U]}}}

\newtheorem{theorem}{Theorem}
\newtheorem{corollary}[theorem]{Corollary}
\newtheorem{proposition}{Proposition}
\newtheorem{definition}{Definition}
\newtheorem{remark}{Remark}

\title{\bfseries A high-order meshless departure-point interpolation operator\\
for semi-Lagrangian lattice Boltzmann methods:\\[0.25em]
\large theory, provenance, and algorithm specification}
\author{Edit17\_CCSLBM project --- internal technical report}
\date{11 August 2026}

\begin{document}
\maketitle

\begin{abstract}
\noindent
This report specifies the theoretical basis for replacing the tensor-product
Lagrange interpolation of an interpolation-supplemented (semi-Lagrangian) lattice
Boltzmann solver with an operator constructed by the Local Anisotropic Basis
Function Method (LABFM) of King, Lind and Nasar~\cite{KLN2020,KL2021}. The central
observation is that the LABFM moment matrix is independent of the operator being
approximated; consequently a single local factorisation yields, at no additional
matrix cost, a family of \emph{off-node interpolation} operators---one per discrete
velocity---in addition to the differential operators for which LABFM was originally
formulated. Polynomial reproduction, constant preservation, the degenerate
zero-displacement limit, the truncation error and its amplification constant, a von
Neumann stability criterion for the streaming step, and invariance under local affine
rescaling are each established with proof. A generalised moving least squares
formulation of the same departure-point operator is specified alongside as a control
arm; it is proved to possess identical polynomial reproduction, so that the order of
accuracy is \emph{not} a discriminator between the two constructions and the choice
between them reduces to a single structural difference in the conditioning of the
local system. Every statement is labelled either with an
explicit equation-level citation to the source publication, with a proof given here,
or as an open item that must be measured before it may be used as an implementation
assumption. One correction is load-bearing: the boundary claim of the 2020 paper, that
incomplete stencils yield one-sided differential operators of the internal order, is
superseded by the 2021 paper, in which the same authors report their tested construction
to be unusable above fourth order and unstable at fourth order, and replace it by a
finite-difference hybrid. This adverse result transfers as a conditioning warning for a
fixed moment matrix, but not as a proof that every differently selected and scaled
one-sided interpolation cloud must fail. The present P6/31 selector therefore accepts
wall clouds only through explicit local hard gates, while global interpolation stability
remains to be measured. A verification ladder is provided whose rungs reproduce
published results, so that an implementation
can be audited against the primary literature at each stage. The literature does not
supply a universal support construction or a three-dimensional critical stencil ratio.
Section~\ref{sec:project-stencil} therefore records, separately and explicitly as a
project design rather than a literature result, the current two-dimensional physical
$P_6$/31-point stencil selector, including its mandatory candidate sets, rank and
condition-number gates, and compact-first repair policy. A future fully three-dimensional
selector remains an open item.
\end{abstract}

\medskip
\begin{center}
\begin{tabularx}{\textwidth}{lX}
\toprule
Label & Meaning \\
\midrule
\cite{KLN2020} & King, Lind \& Nasar, \emph{J.\ Comput.\ Phys.} \textbf{415} (2020)
109549. \textbf{All section and equation numbers refer to the published version}, a copy
of which is held in \texttt{docs/literature/}; preprint numbering is not used. \\
\cite{KL2021} & King \& Lind, \emph{J.\ Comput.\ Phys.} \textbf{449} (2022) 110760.
Same convention. Note that this paper renumbers the original-ABF definition
\cite[Eq.~(19)]{KLN2020} as its own Eq.~(8). \\
\New & Derived in the present report; a proof or an explicit construction is given. \\
\TBD & \textbf{Open item.} Not established by the literature and not proved here.
Must be measured. No such item may be used as an implementation assumption. \\
\bottomrule
\end{tabularx}
\end{center}

\tableofcontents
\newpage

%========================================================================
\section{Scope and position with respect to the literature}

\subsection{Scope}

The target solver discretises the discrete-velocity Boltzmann equation on a
body-fitted curvilinear grid and advances it by an interpolation-supplemented
(semi-Lagrangian) scheme: a collision step evaluated pointwise, followed by transport
along the exact characteristics of the free-streaming operator, with the distribution
function at the departure point recovered by interpolation. In the present solver
that interpolation is a tensor-product seven-point Lagrange interpolation performed in
\emph{computational} space; because the characteristic is a straight line in physical
space but not in computational space, the departure point must be obtained by
numerical integration of the contravariant velocity, which in turn requires the metric
terms of the mapping and a free-stream preservation constraint on their discretisation.

This report specifies an alternative in which the interpolation is performed directly
in \emph{physical} space on a node cloud, using operators constructed by LABFM. The
departure point is then available in closed form, and the metric terms together with
the characteristic integration are eliminated. The collision model, the forcing scheme,
the statistics, the checkpointing and the domain decomposition are unchanged.

The concrete selector presently implemented and audited in
\path{Lagrange_vs_LABFM/2.labfm/} reconstructs in the physical $(y,z)$ plane: it uses
the two-dimensional degree-six space ($p=27$ nonconstant moments), thirty physical
neighbours and a structural centre weight. Section~\ref{sec:project-stencil} is the
normative description of that implementation. Statements below about an 83-dimensional
three-dimensional basis describe the prospective generalisation, not the current
P6/31-point code path.

\begin{center}
\fbox{\parbox{0.95\textwidth}{
\textbf{What is \emph{not} eliminated.} An earlier draft of this report also claimed
that the solver-wide near-wall ghost-layer infrastructure would be removed, on the
strength of the abstract of~\cite{KLN2020} (``At domain boundaries (with incomplete
support) LABFM automatically provides one-sided differences of the same order as the
internal scheme, up to 4th order''). \textbf{That blanket claim is withdrawn.}
\cite[\S 7.1]{KL2021} instead reports failure or instability for its one-sided
differential construction and adopts a finite-difference hybrid. The current LABFM
selector itself reads physical nodes only and clamps its cloud into the domain, but the
existing ghosts remain required by the wall boundary, incumbent ITB/GILBM paths and
statistics contract. Local P6/31 solve gates do not establish global wall stability.
See \S\S\ref{sec:project-stencil}, \ref{sec:boundary} and open item U4.
}}
\end{center}

\subsection{Why the operator route is not pursued}

The alternative of discretising $\partial_t f_q + \bm e_q\!\cdot\!\nabla f_q =
\Omega_q$ directly, using LABFM \emph{differential} operators and an explicit
Runge--Kutta integrator, is not pursued here, for two reasons. First, it has been
realised with radial basis function-generated finite differences~\cite{Maidenberg2026}.
Second, and decisively, the convective operators of LABFM are shown
in~\cite[\S 4.2]{KLN2020} to possess eigenvalues with strictly positive real part on
any node distribution carrying disorder, for every fundamental radial basis function
tested; purely convective problems are therefore unconditionally unstable and require
hyperviscosity~\cite[Eq.~(25)]{KLN2020}. The same conclusion is reached independently
in~\cite[Fig.~10]{Broadley2026}. For a wall-resolved simulation of separated
turbulence, artificial dissipation of the magnitude required would contaminate the
wall shear stress, which is the primary quantity of interest.

The semi-Lagrangian route avoids this failure mode entirely, because transport is
effected by translation along exact characteristics rather than by a discrete
gradient. The relevant stability question becomes that of the interpolation operator,
which is addressed in \S\ref{sec:vn}.

\subsection{Prior art}

Meshless lattice Boltzmann formulations date from~\cite{Musavi2015}. Semi-Lagrangian
streaming with local radial basis function interpolation was introduced
in~\cite{LinWuZhang2019}, and with moving least squares reconstruction
in~\cite{WangYang2025}. The convergence behaviour of the meshless lattice Boltzmann
method has been characterised in~\cite{Strzelczyk2023,Strzelczyk2024}, where a
saturation law is established: the error of the meshless scheme at a given streaming
distance is bounded below by the error of the standard lattice Boltzmann method at
that same streaming distance, so that node refinement ceases to reduce the error once
the interpolation contribution falls below that floor. Separately, high-order
semi-Lagrangian kinetic schemes on body-fitted element meshes have been advanced to
three-dimensional wall-bounded turbulence~\cite{Kramer2017,Wilde2020,Wilde2021,%
Spelten2024}, and interpolation-based lattice Boltzmann methods have been applied to
direct numerical simulation of turbulent channel and duct flows~\cite{Jin2025}.

Against this background, the contribution specified here is narrow and is stated
precisely: (i) LABFM has not previously been coupled to a kinetic scheme; (ii) LABFM
has been formulated only for differential operators, and its extension to off-node
interpolation is given in \S\ref{sec:interp} of this report; and (iii) the polynomial
consistency \emph{realised} here exceeds that reported for any existing meshless
lattice Boltzmann formulation, in three dimensions and for wall-resolved separated
turbulence. Item~(iii) is deliberately a statement about realisation and not about
capability: Theorem~\ref{thm:mls} establishes that the same consistency order is
attainable by moving least squares, so high order on a node cloud is not specific to
LABFM. Claims of novelty beyond these three are not supported by the literature survey
and are not made.

%========================================================================
\section{Notation}

The domain is discretised by a node cloud. Node $i$ is at $\bm r_i$; the difference of
any quantity between nodes is written $(\cdot)_{ji}=(\cdot)_j-(\cdot)_i$ and
$\rji=\bm r_j-\bm r_i$. Following~\cite{KL2021}, $s_i$ denotes the mean node spacing
in the neighbourhood of node $i$ and $h_i$ the stencil length scale; the support radius
is $2h_i$ throughout~\cite[\S 2]{KLN2020}, so that the computational stencil is
$\mathcal N_i=\{\,j : |\rji|\le 2h_i\,\}$ with cardinality $N_i$. The order of
polynomial reproduction is $m$ (denoted $k$ in~\cite{KLN2020} and $m$
in~\cite{KL2021}; the latter convention is used here). The number of spatial
dimensions is $\mu$. Lattice quantities follow the usual D3Q19 convention, with
$\bm e_q$, $q=0,\dots,18$, the discrete velocities and $\Delta t$ the global time step.

\begin{center}
\fbox{\parbox{0.95\textwidth}{
\textbf{Convention: ``order'' is not used as a primary index.} The word is read
differently in the two settings this report has to bridge, and the difference is exactly
one. For a \emph{derivative} operator, reproduction of $\mathbb P_m$ gives $m$-th order
accuracy, and \cite{KLN2020,KL2021} accordingly write ``$m$-th order'' for degree $m$;
their statement that one-sided LABFM is unstable ``at fourth order'' means at $m=4$. For
an \emph{interpolation} operator, reproduction of $\mathbb P_m$ gives truncation
$O(h^{m+1})$, so the same degree is conventionally called $(m{+}1)$-th order. A
seven-point one-dimensional Lagrange interpolant is therefore ``degree six'' under the
first reading and ``seventh order'' under the second, and both descriptions are correct.

\smallskip
\noindent To remove the ambiguity, \textbf{every statement in this report is indexed by
the polynomial degree $m$, with the truncation exponent written explicitly where it
matters}. Ordinal words are used only when quoting a source, and are then tied to
the degree. The incumbent scheme of the target solver is a seven-point tensor Lagrange
interpolant, i.e.\ $m=6$, truncation $O(h^{7})$.
}}
\end{center}

\noindent The basis dimension follows from $m$ and $\mu$ alone. For interpolation in the
streamwise--wall-normal plane the relevant counts are

\begin{center}\small
\begin{tabularx}{\textwidth}{@{}cccccX@{}}
\toprule
\textbf{1-D points} & \textbf{degree $m$} & \textbf{truncation} &
$\dim\mathbb P_m^{2\rm D}$ & \textbf{in-plane neighbours $p_2$} &
\textbf{tensor in-plane points}\\
\midrule
5 & 4 & $O(h^{5})$ & 15 & 14 & $25$\\
6 & 5 & $O(h^{6})$ & 21 & 20 & $36$\\
\textbf{7 (incumbent)} & \textbf{6} & $O(h^{7})$ & \textbf{28} & \textbf{27} &
$\bm{49}$\\
8 & 7 & $O(h^{8})$ & 36 & 35 & $64$\\
\bottomrule
\end{tabularx}
\end{center}

\noindent $p_2=\dim\mathbb P_m^{2\rm D}-1$, the constant being carried structurally by
the increment form and not by a neighbour.

%========================================================================
\section{The local anisotropic basis function method}
\label{sec:labfm}

This section is a review. Every equation carries an explicit citation to its source;
Appendix~\ref{app:prov} collects the correspondence in tabular form.

\subsection{Discrete operators and the moment system}

Define the vector of monomials and the vector of partial derivative operators
\cite[Eqs.~(1),(2)]{KLN2020}, \cite[Eqs.~(1),(2)]{KL2021}
\begin{align}
\bm X_{ji}&=\Bigl[x_{ji},\ y_{ji},\ \tfrac{x_{ji}^{2}}{2},\ x_{ji}y_{ji},\
\tfrac{y_{ji}^{2}}{2},\ \tfrac{x_{ji}^{3}}{6},\ \dots,\ \tfrac{y_{ji}^{m}}{m!}
\Bigr]^{\mathsf T},
\label{eq:X}\\[0.2em]
\bm D(\cdot)&=\bigl[\partial_x,\ \partial_y,\ \partial_{xx},\ \partial_{xy},\
\partial_{yy},\ \partial_{xxx},\ \dots,\ \partial_{y^m}\bigr]^{\mathsf T}(\cdot).
\label{eq:D}
\end{align}
It is emphasised, as in~\cite{KLN2020}, that $\bm X_{ji}\neq\bm X_j-\bm X_i$: the
entries of $\bm X_{ji}$ are monomials in the separation vector $\rji$. With this
convention the multidimensional Taylor expansion takes the compact form
\cite[Eq.~(3)]{KLN2020}
\begin{equation}
(\cdot)_j=(\cdot)_i+\bm X_{ji}\cdot\bm D(\cdot)\big|_i .
\label{eq:taylor}
\end{equation}

A general discrete operator is written \cite[Eq.~(4)]{KLN2020},
\cite[Eq.~(3)]{KL2021}
\begin{equation}
\mathcal L^{d}_i(\cdot)=\sum_{j\in\mathcal N_i}(\cdot)_{ji}\,w^{d}_{ji}
\ \approx\ \bm C^{d}\cdot\bm D(\cdot)\big|_i ,
\label{eq:Lop}
\end{equation}
in which $\bm C^{d}$ selects the derivative combination to be approximated
\cite[Eq.~(5)]{KLN2020}, \cite[Eq.~(4)]{KL2021}; for the gradient components and the
Laplacian, $\bm C^{x}=[1,0,0,\dots]^{\mathsf T}$, $\bm C^{y}=[0,1,0,\dots]^{\mathsf T}$
and $\bm C^{L}=[0,0,1,0,1,0,\dots]^{\mathsf T}$.

The weights are expanded in a set of anisotropic basis functions (ABFs)
$\bm W_{ji}=\bm W(\rji/h_i)$ \cite[Eq.~(10)]{KLN2020}, \cite[Eq.~(5)]{KL2021},
\begin{equation}
w^{d}_{ji}=\bm W_{ji}\cdot\bm\Psi^{d}_i ,
\label{eq:wPsi}
\end{equation}
which upon substitution into the moment vector $\bm B^d_i=\sum_j\bm X_{ji}w^d_{ji}$
\cite[Eq.~(8)]{KLN2020} yields the local linear system
\cite[Eqs.~(13),(14)]{KLN2020}, \cite[Eqs.~(6),(7)]{KL2021}
\begin{equation}
\bm M_i=\sum_{j\in\mathcal N_i}\bm X_{ji}\otimes\bm W_{ji},
\qquad
\bm M_i\,\bm\Psi^{d}_i=\bm C^{d}.
\label{eq:moments}
\end{equation}

\begin{remark}[Operator independence of the moment matrix]
\label{rem:shared}
The matrix $\bm M_i$ does not depend on $d$. This is stated explicitly
in~\cite[\S 3.1]{KLN2020}: ``The matrix $M_i$ is the same for all operators
(different $d$), and so we calculate each $M_i$ once, then solve (14) for each $d$ to
obtain the required operators.'' The construction of \S\ref{sec:interp} rests entirely
on this property: it replaces $\bm C^d$ by a family of right-hand sides indexed by the
discrete velocity, so that one factorisation per node serves all streaming directions
as well as every post-processing differential operator.
\end{remark}

\begin{remark}[Distributional reading: what $\bm C$ is; \New]
\label{rem:distribution}
Associate with the weights the discrete measure
\begin{equation}
\mu^{d}_i\;\equiv\;\sum_{j\in\mathcal N_i}w^{d}_{ji}
\bigl(\delta_{\bm r_j}-\delta_{\bm r_i}\bigr),
\qquad\text{so that}\qquad
\bigl\langle\mu^{d}_i,\phi\bigr\rangle=\mathcal L^{d}_i(\phi).
\label{eq:measure}
\end{equation}
The moment conditions \eqref{eq:moments} then read: \emph{$\mu^d_i$ agrees with a
prescribed target distribution when paired with every element of the test space
$\mathbb P_m$}, and $\bm C^d$ is precisely the vector of moments of that target,
$C^{d,\beta}=\langle\text{target},X^\beta\rangle$. Explicitly,
\begin{center}
\begin{tabular}{ll}
\toprule
Target distribution & Resulting $\bm C$ \\
\midrule
$-\partial_x\delta_{\bm r_i}$ & $\bm C^{x}=[1,0,0,\dots]^{\mathsf T}$ \\
$\nabla^{2}\delta_{\bm r_i}$ & $\bm C^{L}=[0,0,1,0,1,0,\dots]^{\mathsf T}$ \\
$\delta_{\bm r_i+\bm\delta_q}-\delta_{\bm r_i}$ & $\bm C^{q}=\bm X(\bm\delta_q)$
\quad(\S\ref{sec:interp}) \\
\bottomrule
\end{tabular}
\end{center}
The last line follows from
$\langle\delta_{\bm r_i+\bm\delta_q}-\delta_{\bm r_i},X^\beta(\cdot-\bm r_i)\rangle
=X^\beta(\bm\delta_q)-X^\beta(\bm 0)=X^\beta(\bm\delta_q)$, and incidentally explains
why $\bm X$ carries no constant term: the difference form annihilates constants on both
sides. This reading unifies the operators of~\cite{KLN2020,KL2021} with the one
introduced in \S\ref{sec:interp} --- they differ only in the target distribution.
\end{remark}

\subsection{Order of accuracy and basis dimension}

When \eqref{eq:moments} is solved with $\bm M_i$ of size $p\times p$, exact
reproduction of polynomials of total degree $m$ is enforced and the leading truncation
error of $\mathcal L^d_i$ is $O\!\left(h_i^{\,m+1-l}\right)$, where $l$ is the order of
the derivative approximated \cite[\S 3.1]{KLN2020}, \cite[\S 2]{KL2021}.
Accordingly the gradient converges at order $m$, the Laplacian at order $m-1$, and
--- of direct relevance here --- an operator with $l=0$, that is an interpolation,
converges at order $m+1$.

The basis dimension is \cite[Eqs.~(15),(16)]{KLN2020}
\begin{equation}
p_{\mathrm{2D}}=\frac{m^{2}+3m}{2},
\qquad
p_{\mathrm{3D}}=\sum_{q=2}^{m+1}\frac{q(q+1)}{2}=\binom{m+3}{3}-1 .
\label{eq:p}
\end{equation}

\begin{center}
\begin{tabular}{lrrrrrr}
\toprule
$m$ & 2 & 4 & 5 & 6 & 8 & 10\\
\midrule
$p_{\mathrm{2D}}$ & 5 & 14 & 20 & 27 & 44 & 65\\
$p_{\mathrm{3D}}$ & 9 & 34 & 55 & \textbf{83} & 164 & 285\\
interpolation order $m+1$ & 3 & 5 & 6 & \textbf{7} & 9 & 11\\
\bottomrule
\end{tabular}
\end{center}

\noindent The three-dimensional values reproduce the sequence $3,9,19,34,55$ quoted
in \cite[Eq.~(16)]{KLN2020} for $m=1,\dots,5$. The polynomial design point is
$m=6$ (truncation $O(h^{7})$), matching the degree of the seven-point Lagrange
interpolation it replaces. The current in-plane implementation has
$p_{\mathrm{2D}}=27$ and is specified in \S\ref{sec:project-stencil}; a genuinely
three-dimensional implementation would instead have $p_{\mathrm{3D}}=83$ and is not
specified by the present source files.

\subsection{Structure of the weight system: what the basis functions decide}
\label{sec:structure}

The material of this subsection is not stated in this form in~\cite{KLN2020,KL2021},
but it is the shortest route to understanding both the role of the ABFs and the
relationship to the control formulation of \S\ref{sec:mls}. \New

Collect the stencil data into matrices: let $\bm X\in\mathbb R^{N_i\times p}$ have rows
$\widetilde{\bm X}_{ji}^{\mathsf T}$ and $\bm W\in\mathbb R^{N_i\times p}$ have rows
$\bm W_{ji}^{\mathsf T}$, and let $\bm w\in\mathbb R^{N_i}$ collect the weights. The
moment conditions \eqref{eq:moments} are
\begin{equation}
\bm X^{\mathsf T}\bm w=\bm C
\qquad\text{---\ $p$ equations in $N_i$ unknowns.}
\label{eq:underdet}
\end{equation}
For $N_i>p$ this system is underdetermined: its solution set is an affine subspace of
dimension $N_i-p$, and \emph{every member reproduces $\mathbb P_m$ exactly}. A
selection rule is therefore mandatory. The ABF ansatz \eqref{eq:wPsi} supplies one by
restricting $\bm w$ to the column space of $\bm W$, which renders the system square,
and $\bm M_i=\bm X^{\mathsf T}\bm W$ is exactly that restriction. Writing the rule with
a general \emph{selector} $\bm S\in\mathbb R^{N_i\times p}$,
\begin{equation}
\boxed{\ \bm w=\bm S\bigl(\bm X^{\mathsf T}\bm S\bigr)^{-1}\bm C\ }
\qquad
\begin{cases}
\bm S=\bm W & \text{LABFM},\\[0.15em]
\bm S=\bm\Theta\bm Q & \text{generalised moving least squares (\S\ref{sec:mls})}.
\end{cases}
\label{eq:selector}
\end{equation}

\begin{proposition}[Degenerate limit $N_i=p$; \New]
\label{prop:degenerate}
If $N_i=p$ and both $\bm X$ and $\bm S$ are invertible, then
$\bm w=\bm X^{-\mathsf T}\bm C$, independent of $\bm S$. The weights are then those of
plain polynomial interpolation of degree $m$; in one dimension with $m$ neighbours plus
the centre they are the Lagrange weights on $m+1$ points.
\end{proposition}
\begin{proof}
$\bm w=\bm S(\bm X^{\mathsf T}\bm S)^{-1}\bm C=\bm S\bm S^{-1}\bm X^{-\mathsf T}\bm C
=\bm X^{-\mathsf T}\bm C$. Uniqueness of the solution of the square system
\eqref{eq:underdet} identifies it with the unique degree-$m$ polynomial reproduction
rule on those $N_i+1$ points.
\end{proof}

\begin{remark}[What the surplus degrees of freedom buy]
\label{rem:surplus}
Proposition~\ref{prop:degenerate} shows that at $N_i=p$ the basis functions are
irrelevant: they cancel. They begin to act only for $N_i>p$, and by
\eqref{eq:underdet} every admissible choice has \emph{identical} polynomial
consistency. The surplus $N_i-p$ degrees of freedom therefore cannot change the order of accuracy;
what they do change is the coefficient of the leading error term, together with the
conditioning and stability of the resulting operator (Appendix~\ref{app:tutorial}
exhibits both statements explicitly). This is consistent with the requirement
$N_{\rm crit}>N_{\rm poly}$ reported in \cite[\S 3.2]{KLN2020} --- for instance $37$
against $28$ at $m=6$ in two dimensions --- and it is the precise sense in which LABFM
and moving least squares differ: same constraint set, different selector.
\end{remark}

\begin{proposition}[The expansion origin is a gauge; \New]
\label{prop:gauge}
The constraint \eqref{eq:underdet} is invariant under a change of the point about which
the monomials $\bm X$ and the right-hand side $\bm C$ are formed. Consequently, with
$\bm S$ and $\mathcal N_i$ held fixed, the weights $\bm w$ do not depend on that choice.
\end{proposition}
\begin{proof}
The constraint is equivalent to the origin-free statement ``$\sum_j w_{ji}(P_j-P_i)=
\langle\text{target},P\rangle$ for all $P\in\mathbb P_m$''. Monomials about any point
form a basis of $\mathbb P_m$ (the constant contributing the trivial identity $0=0$ in
the difference form), so different choices produce different but equivalent bases for
the same $p$ conditions, hence the same solution set; intersecting with the fixed
subspace $\operatorname{span}(\bm S)$ gives the same $\bm w$.
\end{proof}

\begin{remark}[Three distinct ``origins'', only one of which is free]
\label{rem:origins}
Formulae in~\cite{KLN2020} place three things at $\bm r_i$, and they are commonly
conflated: (i) the monomial/Taylor origin, which is a gauge by
Proposition~\ref{prop:gauge}; (ii) the ABF argument origin, which is \emph{not} a gauge,
since it defines $\operatorname{span}(\bm S)$ and hence which admissible solution is
selected, together with $\kappa$, $\Lambda$ and $|G_q|$; and (iii) the stencil centre,
which is not a gauge either, since it determines $\mathcal N_i$. Gauge freedom in (i) is
nevertheless not licence: with the origin far from the stencil the entries of $\bm X$
span many orders of magnitude and $\bm M_i$ becomes numerically intractable although
formally equivalent. Choosing $\bm r_i$, together with the scaling of
\S\ref{sec:affine}, is the well-conditioned gauge. Proposition~\ref{prop:gauge} is used
as a verification test in rung L5f of \S\ref{sec:ladder}. The alternative of centring
(ii) and (iii) on the departure point is a genuine method variant, adopted for example
in the meshless lattice Boltzmann scheme of~\cite{Strzelczyk2023}; it costs one moment
matrix per discrete velocity.
\end{remark}

\subsection{Construction of the anisotropic basis functions}
\label{sec:abf}

Two constructions appear in the literature.

\paragraph{Original ABFs.} In~\cite{KLN2020} the ABFs are generated as partial
derivatives of a fundamental radial basis function $W^{0}$,
\begin{equation}
\bm W_{ji}=\bm D\,W^{0}(\rji),
\label{eq:origABF}
\end{equation}
\cite[Eq.~(19)]{KLN2020}, with four choices of $W^{0}$ examined (conic, quadratic,
Wendland C6, Gaussian; explicit forms in \cite[Appendix A]{KLN2020}). The fundamental
functions are normalised so that $W^{0}(2)=0$ (Gaussian excepted) and
$2\pi\int_0^{2}qW^{0}(q)\,\mathrm dq=1$ \cite[Eq.~(A.1)]{KLN2020}; the convex quadratic
is $W_Q(q)=\tfrac{3}{16\pi}(q-2)^2$ on $0\le q\le2$ \cite[Eq.~(A.3)]{KLN2020}, designed
so that $\mathrm dW_Q/\mathrm dq|_{q=2}=0$, which~\cite{KLN2020} identifies as apparently
necessary for stability of the convective operator on a uniform distribution.

\begin{remark}[Implementation caveat required to reproduce~\cite{KLN2020}; \New]
\label{rem:truncation}
The full expansion \cite[Eq.~(A.2)]{KLN2020} of $\bm D W^{0}$ contains terms in
$\mathrm d^2W^{0}/\mathrm dr^2$ and higher radial derivatives, which vanish only for the
conic $W^{0}$. \cite[Appendix A]{KLN2020} states that in practice these terms are
\emph{omitted} for all $m\ge4$ (that is, from the tenth element of $\bm W_{ji}$ onward)
\emph{regardless of whether they vanish}, for two stated reasons: computational
simplicity, and greater resilience to node-distribution disorder. Any attempt to
reproduce the published two-dimensional convergence and critical-ratio results ---
rungs L1 and L2 of \S\ref{sec:ladder} --- must apply the same truncation, or the
numbers will not match. This caveat does not affect the Hermite construction, which is
what the production path uses.
\end{remark}

\paragraph{Hermite ABFs.} In~\cite{KL2021} the ABFs are instead built from an
orthogonal polynomial basis. Index the basis by the multi-index
$\alpha$ of the corresponding monomial: if the $\alpha$-th entry of $\bm X_{ji}$ is
proportional to $x_{ji}^{a}y_{ji}^{b}$, then \cite[Eqs.~(10),(11)]{KL2021}
\begin{equation}
W^{\alpha}_{ji}=\frac{\psi\!\left(|\rji|/h_i\right)}{\sqrt{2^{\,a+b}}}\,
H_a\!\left(\frac{x_{ji}}{h_i\sqrt2}\right)
H_b\!\left(\frac{y_{ji}}{h_i\sqrt2}\right),
\label{eq:hermiteABF}
\end{equation}
where $H_a$ is the physicists' Hermite polynomial of degree $a$ and $\psi$ is a radial
basis function (Wendland C2 by default; \cite{KL2021} states that other choices,
including $\psi\equiv1$, are admissible).

\begin{remark}[The roles of the two subscripts, and the shape of $\bm W$]
\label{rem:indices}
$W^{\alpha}_{ji}$ carries three indices: $\alpha$ selects the basis function, $j$ the
neighbour at which it is evaluated, and $i$ the centre. The object
$\bm W_{ji}=[W^{1}_{ji},\dots,W^{p}_{ji}]^{\mathsf T}$ appearing in \eqref{eq:wPsi} is
therefore, for one \emph{fixed} neighbour, the vector of \emph{all $p$} basis functions
evaluated at the single separation $\rji$. Every neighbour carries every basis function;
the basis is not distributed among the neighbours.

Two consequences are worth stating, because the alternative reading is a common one.
First, dimensionally: only if $\bm W_{ji}$ is a $p$-vector is
$\bm X_{ji}\otimes\bm W_{ji}$ a $p\times p$ matrix and $\bm M_i$ a square system; were
each neighbour to carry a single basis function, the sum in \eqref{eq:moments} would
return a vector and the construction would not close. Equivalently, in the notation of
\S\ref{sec:structure}, $\bm W$ is an $N_i\times p$ matrix with neighbours indexing rows
and basis functions indexing columns --- the same shape as $\bm X$, and the same shape as
a design matrix in least squares.

Second, by contrast: radial basis function methods \emph{do} index the basis by node,
placing one function centred at each neighbour, which yields an $N_i\times N_i$ system
whose size grows with the stencil. In LABFM the basis is indexed by differential order
instead, so enlarging the stencil adds equations to an overdetermined set rather than
unknowns, and the local system remains $p\times p$ whatever $N_i$ is. This is the origin
of the $O(p^{3})$ rather than $O(N_i^{3})$ factorisation cost, and of the surplus
degrees of freedom discussed in Remark~\ref{rem:surplus}.
\end{remark}

\begin{center}
\fbox{\parbox{0.95\textwidth}{
\textbf{Decisive result for the three-dimensional case.} \cite[\S 3]{KL2021} reports:
``An example of where this results in the method breaking down is when LABFM is
extended to three-dimensions, with $m=4$. Here $M$ has size $n=34$, and rank $n-1=33$,
\emph{regardless of the node distribution}.''

\smallskip
The original ABFs \eqref{eq:origABF} therefore render $\bm M$ rank deficient in three
dimensions for $m\ge4$, independently of the node cloud. The present application is
three-dimensional with target $m=6$; the Hermite construction \eqref{eq:hermiteABF} is
consequently \emph{mandatory}. The original construction retains value only as a
two-dimensional regression target (rungs L1--L2 of \S\ref{sec:ladder}) and as the
subject of the negative test L3, and must not appear on any three-dimensional
production path.
}}
\end{center}

\subsection{Row preconditioning}

Define the vector of inverse powers of $h_i$ \cite[Eq.~(9)]{KL2021}
\begin{equation}
\bm H_i=\bigl[h_i^{-1},h_i^{-1},h_i^{-2},h_i^{-2},h_i^{-2},h_i^{-3},\dots,
h_i^{-m}\bigr]^{\mathsf T},
\qquad
\widetilde{\bm X}_{ji}=\bm X_{ji}\!\cdot\!\bm H_i,
\quad
\widetilde{\bm C}^{d}=\bm C^{d}\!\cdot\!\bm H_i,
\label{eq:precond}
\end{equation}
the products being taken entrywise. Solving \eqref{eq:moments} with
$\widetilde{\bm X}$ and $\widetilde{\bm C}$ in place of $\bm X$ and $\bm C$ leaves the
exact solution unchanged while improving the condition number of $\bm M$ by several
orders of magnitude~\cite[\S 2]{KL2021}. This preconditioning is applied
unconditionally in the present specification.

\subsection{Sampling requirements and stencil sizing}
\label{sec:sampling}

The construction requires that the discrete sums forming $\bm M_i$ approximate the
corresponding continuous integrals~\cite[Eq.~(22)]{KLN2020}. This imposes a lower bound
on $h/s$, made precise in~\cite[\S 3.2]{KLN2020} by an angular sampling argument:
expressed in local polar coordinates, the order-$m$ ABFs contain terms in $\cos m\theta$
and $\sin m\theta$, so that if the support is divided into $2m$ equal angular sectors,
each sector must contain at least one node. Below the resulting critical ratio, some
eigenvalues of $\bm M_i$ fall to machine-precision zero and the condition number
increases by roughly ten orders of magnitude~\cite[\S 3.2]{KLN2020}.

\begin{remark}[Normalisation of the sampling condition; \New]
\label{rem:innerprod}
The inner product of \cite[Eq.~(18)]{KLN2020} is a \emph{mean} over the support,
$\langle f,g\rangle=V_i^{-1}\int_{V_i}fg\,\mathrm dV_j$, so that
\cite[Eq.~(22)]{KLN2020} as printed, $\sum_j X^{a}_{ji}W^{b}_{ji}\approx\langle
X^{a},W^{b}\rangle$, is short of a factor $N_i$. With node number density $s^{-\mu}$ the
operative relation is
\begin{equation}
\sum_{j\in\mathcal N_i}X^{a}_{ji}W^{b}_{ji}
\;\approx\;\frac{1}{s_i^{\mu}}\int_{V_i}X^{a}W^{b}\,\mathrm dV
\;=\;N_i\bigl\langle X^{a},W^{b}\bigr\rangle .
\label{eq:sampling-norm}
\end{equation}
The distinction is immaterial to the argument of~\cite{KLN2020} but must be respected
when a continuum reference matrix is computed as a diagnostic (\S\ref{sec:parity}).
\end{remark}

\paragraph{Published two-dimensional values.} These belong to \emph{two different}
ABF constructions and must not be conflated:
\begin{center}
\begin{tabular}{llrrrrr}
\toprule
Quantity & ABF used & $m{=}2$ & 4 & 6 & 8 & 10\\
\midrule
$(h/s)_{\rm crit}$, \emph{critical} \cite[\S 3.2]{KLN2020} & original & 0.75 & 1.15
& 1.60 & 2.10 & ---\\
$N_{\rm crit}$ \cite[\S 3.2]{KLN2020} & original & 8 & 21 & 37 & 57 & ---\\
$N_{\rm poly}=p_{\rm 2D}+1$ & --- & 6 & 15 & 28 & 45 & 66\\
$h/s$, \emph{recommended} \cite[\S 5]{KL2021} & \textbf{Hermite} & 1.2 & 1.4
& \textbf{1.8} & 2.3 & 2.8\\
\bottomrule
\end{tabular}
\end{center}
\noindent The stencil optimisation of~\cite[\S 4]{KL2021} reduces the mean ratio for
$m=6$ from $2.2$ to approximately $1.5$, lowering the neighbour count by $17.5\,\%$.

\begin{center}
\fbox{\parbox{0.95\textwidth}{
\TBD\ \textbf{Open item U8.} No \emph{critical} ratio is published for the Hermite
ABFs, in any dimension. The critical values above were measured with the original
ABFs, which the production path does not use; the Hermite values quoted
by~\cite{KL2021} are recommended operating points carrying an unquantified margin.
Code must therefore not name the original-ABF critical values as if they applied to the
production construction.
}}
\end{center}

The neighbour count follows from the support geometry. In two dimensions
$N_i\approx4\pi(h_i/s_i)^{2}$~\cite[\S 2]{KL2021}; the same argument in three
dimensions gives \New
\begin{equation}
N_i^{\rm 3D}\approx\frac{4}{3}\pi\frac{(2h_i)^{3}}{s_i^{3}}
=\frac{32\pi}{3}\left(\frac{h_i}{s_i}\right)^{3}
\approx33.51\left(\frac{h_i}{s_i}\right)^{3}.
\label{eq:N3D}
\end{equation}

\begin{center}
\fbox{\parbox{0.95\textwidth}{
\TBD\ \textbf{Open item U1.} The critical ratio $(h/s)^{\rm 3D}_{\rm crit}$ is not
available. The values tabulated above are two-dimensional. The three-dimensional
sampling condition is a solid-angle rather than a plane-angle condition and is
necessarily more demanding, but by how much is not established;~\cite{KL2021} states
explicitly that three-dimensional simulations were reserved for a later study. Because
$N_i$ scales as $(h/s)^{3}$, this single unknown controls any cost model built on
the neighbour count.
}}
\end{center}

\subsection{Current project-specific \texorpdfstring{$P_6$/31-point}{P6/31-point}
stencil selector}
\label{sec:project-stencil}

This subsection is an implementation specification, not a rule taken from
\cite{KLN2020,KL2021}. Its normative sources are
\path{Lagrange_vs_LABFM/2.labfm/labfm_stencil.h} and
\path{Lagrange_vs_LABFM/2.labfm/labfm_solve.h}, read through the active call chain
\texttt{b\_weights\_case2} $\rightarrow$
\texttt{labfm\_select\_physical\_stencil} $\rightarrow$
\texttt{labfm\_candidate\_pool\_build}. The separately named
\texttt{labfm\_select\_physical\_stencil\_legacy\_rrqr} has no caller in this path and
is retained only for validation; it does not define the current selector. The Taylor,
Hermite and Wendland functions called by this path are defined in
\path{Lagrange_vs_LABFM/2.labfm/labfm_hermite.h}.

\begin{center}
\fbox{\parbox{0.95\textwidth}{
\textbf{Project rule (a posteriori certificate, not an a priori topology theorem).}
The following five stages map one-to-one to
\S\S\ref{sec:project-point-sets}--\ref{sec:project-final-solve}:
\begin{enumerate}[label=\arabic*.,leftmargin=2.0em]
\item \textbf{Point sets.} Construct a deterministic physical-node candidate pool and
      distinguish points protected in that pool from points selected in the final
      31-point footprint.
\item \textbf{Compact tier.} First try the fixed seven-layer index-space pyramid. Its
      actual physical shape may be warped, and the template is not assumed to be
      nonsingular before it is tested.
\item \textbf{Admissibility certificate.} Assemble the matrices from the actual physical
      coordinates. Accept a trial only when its support, numerical rank, reciprocal
      conditioning and direction-independent weight-amplification gates all pass.
\item \textbf{Repair tier.} Only if every compact-tier trial fails, generate replacement
      clouds, apply the same admissibility certificate, and compare the surviving clouds
      by a degree-seven error proxy.
\item \textbf{Final solve and export.} For each lattice direction, solve for its weights
      and pass independent solver and backward-error checks. Failed data are never
      exported to the time-step path.
\end{enumerate}
Thus geometry proposes a stencil, algebra certifies it, and only a certified stencil may
be used. Conditioning is an acceptance gate, not the optimisation objective.
}}
\end{center}

\subsubsection{From grid indices to three point sets}
\label{sec:project-point-sets}

\paragraph{Purpose of this stage.}
This stage answers three different questions: which physical nodes are enumerated,
which of them are guaranteed a place in the finite candidate pool, and which 30
neighbours finally accompany the centre. These three sets must not be conflated.

\paragraph{Geometric notation.}
The code stores three ghost layers around the physical array. Let $(j,k)$ be the stored
array indices of the arrival node at which a stencil is being built. Its physical-grid
indices are
\begin{equation}
 c_j=j-3,\qquad c_k=k-3.
\label{eq:project-centre-index}
\end{equation}
Here $N_y$ and $N_z$ are, symmetrically, the numbers of distinct physical-grid nodes in
the streamwise and wall-normal directions. Their canonical physical indices are
$0,\ldots,N_y-1$ and $0,\ldots,N_z-1$, respectively. The streamwise coordinate file
contains one additional record at index $N_y$ that repeats index zero at the other end
of the period. After retaining this repeated endpoint and adding three ghosts on each
side, the stored extents are
\begin{equation}
 N_{y6}=(N_y+1)+6=N_y+7,\qquad N_{z6}=N_z+6.
\label{eq:project-stored-extents}
\end{equation}
Thus the implementation variables satisfy
$\texttt{NI}=N_y+1$ (the solver macro \texttt{NY} has the same convention),
$\texttt{NJ}=N_z$,
$\texttt{NY6}=N_{y6}$, and $\texttt{NZ6}=N_{z6}$; $N_{z6}$ is the stored
wall-normal stride.

For one enumerated candidate, the following symbols are used:
\begin{center}\small
\begin{tabularx}{\textwidth}{@{}P{2.6cm}X@{}}
\toprule
\textbf{Symbol} & \textbf{Meaning} \\
\midrule
$d_j,d_k$ & Integer offsets from $(c_j,c_k)$ in structured index space. They determine
the candidate through $c_j+d_j=j_p+\sigma_p N_y$ and $k_p=c_k+d_k$; $d_j$ is kept
unwrapped so that its periodic image remains known. \\
$j_p,k_p$ & Canonical physical-grid indices of the candidate, with
$0\le j_p<N_y$ and $0\le k_p<N_z$. \\
$\iota_p$ & Canonical physical-node label, stored as \texttt{physical\_id}; it uniquely
identifies $(j_p,k_p)$ but does not distinguish its periodic images. \\
$\sigma_p$ & Integer periodic-image number, stored as \texttt{image\_shift} and uniquely
defined by $c_j+d_j=j_p+\sigma_p N_y$. It identifies which periodic image gives a continuous
streamwise displacement. \\
$L_y$ & Physical streamwise period, stored as \texttt{period\_y}. \\
$Y,Z$ & Physical-coordinate arrays; $Y(j_p,k_p)$ and $Z(j_p,k_p)$ are their values at
the canonical physical node. \\
$\Delta Y_p,\Delta Z_p$ & Physical-coordinate displacement from the arrival node to the
chosen periodic image of the candidate. \\
\bottomrule
\end{tabularx}
\end{center}

If $(Y_0,Z_0)$ denotes the physical coordinate of the arrival node, then
\begin{equation}
 \Delta Y_p=Y(j_p,k_p)+\sigma_pL_y-Y_0,\qquad
 \Delta Z_p=Z(j_p,k_p)-Z_0.
\label{eq:project-physical-displacement}
\end{equation}
The canonical physical-node label used for de-duplication is
\begin{equation}
 \iota_p=j_pN_z+k_p,\qquad 0\le\iota_p<N_y N_z.
\label{eq:project-physical-id}
\end{equation}
The code stores $\iota_p$ in \texttt{physical\_id}. It is not the same as the flattened
storage address $\mathrm{storage\_id}=(j_p+3)N_{z6}+(k_p+3)$, where $N_{z6}$ is the
stored wall-normal stride including ghosts. Two periodic images can have the same
$\iota_p$ and canonical storage address but different values of $\sigma_p$, and hence
different continuous coordinates in \eqref{eq:project-physical-displacement}.

\paragraph{Set 1: all enumerated physical candidates.}
The search is written symmetrically in terms of the two integer offsets:
\begin{equation}
 -16\le d_j\le16,\qquad
 \max(-16,-c_k)\le d_k\le\min(16,N_z-1-c_k).
\label{eq:enum-set}
\end{equation}
For every pair $(d_j,d_k)$, the canonical physical indices are obtained from
$c_j+d_j=j_p+\sigma_p N_y$ and $k_p=c_k+d_k$. The first relation wraps the periodic
streamwise index while retaining its image shift; the bounds on $d_k$ ensure directly
that $0\le k_p<N_z$, so no wall ghost can enter the set. The centre
$(d_j,d_k)=(0,0)$ is skipped. Away from both walls the set therefore contains
$33\mathbin{\times}33-1=1088$ candidates; wall clipping reduces this number. The value
16 is a finite structured-index search budget; it is not the Wendland support radius
and is not a physical-distance cutoff.

To order the candidates, the code first estimates the two local physical spacings
\begin{equation}
\begin{aligned}
 s_y&=\tfrac12\bigl(|Y^{\rm st}_{j+1,k}-Y^{\rm st}_{j,k}|
                   +|Y^{\rm st}_{j,k}-Y^{\rm st}_{j-1,k}|\bigr),\\
 s_z&=\tfrac12\bigl(|Z^{\rm st}_{j,k+1}-Z^{\rm st}_{j,k}|
                   +|Z^{\rm st}_{j,k}-Z^{\rm st}_{j,k-1}|\bigr).
\end{aligned}
\label{eq:project-local-spacing}
\end{equation}
The superscript $\rm st$ indicates that these four differences use stored-array indices;
the three ghost layers make all four accesses valid even at a physical wall.
For each candidate it then defines
\begin{equation}
 \rho_p=\sqrt{(\Delta Y_p/s_y)^2+(\Delta Z_p/s_z)^2},\qquad
 \theta_p=\operatorname{atan2}(\Delta Z_p,\Delta Y_p)\in[-\pi,\pi].
\label{eq:rho-order}
\end{equation}
Here $\rho_p$ is a dimensionless local-grid distance, whereas $\theta_p$ is used only to
break equal-$\rho_p$ ties. The enumerated array is sorted lexicographically by
\begin{equation}
 (\rho_p,\theta_p,\iota_p,\sigma_p).
\label{eq:project-candidate-order}
\end{equation}
Floating-point keys are compared directly; there is no epsilon-equality rule.

\paragraph{Set 2: the candidate pool of at most 160 points.}
\noindent\textbf{Key point:} For an interior arrival node on the audited grids, Set 2 is
not the 160 globally nearest points in Set 1: it first retains the 48 neighbours in a
protected seven-by-seven index-space window, then retains the 112 nearest unique
physical nodes among the other 1040 candidates, where ``nearest'' is determined first
by $\rho_p$ and the remaining keys only break exact ties. Thus
\begin{equation}
 \underbrace{1088}_{\text{Set 1}}
 =\underbrace{48}_{\text{protected}}+\underbrace{1040}_{\text{other candidates}},
 \qquad
 \underbrace{160}_{\text{Set 2}}
 =\underbrace{48}_{\text{protected}}+\underbrace{112}_{\text{nearest retained}}.
\label{eq:set-two-counts}
\end{equation}

\emph{Pass 0: retain the 48 protected neighbours.}
The protected index-space window is
\begin{equation}
\begin{aligned}
 k_{\rm lo}={}&\min\!\bigl(\max(c_k-3,0),N_z-7\bigr),\\
 \mathcal W_{\rm prot}={}&\{(d_j,d_k):-3\le d_j\le3,\;
 k_{\rm lo}-c_k\le d_k\le k_{\rm lo}+6-c_k\}
 \setminus\{(0,0)\}.
\end{aligned}
\label{eq:mandatory-window}
\end{equation}
Here $k_{\rm lo}$ is the first physical $k$-index of the protected seven-row block.
Away from a wall, $k_{\rm lo}=c_k-3$ and hence $-3\le d_k\le3$. At the lower wall
$c_k=0$, the range becomes $0\le d_k\le6$; at the upper wall $c_k=N_z-1$, it becomes
$-6\le d_k\le0$. Thus the same seven-by-seven offset window is shifted inward near a
wall without admitting ghosts. Removing $(0,0)$ excludes the centre, leaving
$7\mathbin{\times}7-1=48$ protected neighbours on the audited grids. These 48 points
enter Set 2 even if some unprotected candidates have smaller $\rho_p$.

\emph{Pass 1: fill the remaining 112 slots by normalised physical distance.}
For an interior node, Pass 1 scans the other $1088-48=1040$ candidates in the order
$ (\rho_p,\theta_p,\iota_p,\sigma_p)$ and inserts the first 112 unique physical nodes.
Before each insertion, the pool is searched for the same $\iota_p$
(\texttt{physical\_id} in the code); a later periodic image of an already retained node
is discarded. Equivalently, the complete priority key is
\[
 (\mathrm{pass},\rho_p,\theta_p,\iota_p,\sigma_p),
\]
with protected-window points assigned to Pass 0. The four candidate keys do not receive
separate point quotas: $\rho_p$ supplies the primary ordering, while $\theta_p$,
$\iota_p$, and $\sigma_p$ are consulted successively only when all preceding keys are
exactly equal.

Near a wall, Set 1 contains fewer than 1088 candidates, so the pool of unprotected
candidates is smaller than 1040; the algorithm itself is unchanged. On the audited
PH/nonPH grids, enough unique candidates still exist, so Pass 0 retains 48 protected
points and Pass 1 still fills the remaining $160-48=112$ slots. If Set 1 were exhausted
earlier, Set 2 would simply contain fewer than 160 points. If this header is reused with
$N_y\le32$, an explicit guard comparing $\iota_p$ with the centre's canonical-node label
must be added before making the same uniqueness claim; the present input check alone
allows $\texttt{NI}=N_y+1\ge8$, that is, $N_y\ge7$.

The protected window is needed because, on a strongly stretched grid, the first 160
points in pure $\rho_p$ order can occupy at most six nearly parallel wall-normal layers.
Such a set can make the degree-six sampling problem algebraically deficient. Protection
keeps a seven-layer construction available to both the compact and repair tiers.
Protection means \emph{must enter the pool}, not \emph{must enter the final stencil}.

\paragraph{Set 3: the final 31-point footprint.}
The final footprint contains the arrival node, labelled $n=0$, and exactly 30 physical
neighbours, labelled $n=1,\ldots,30$. Only the 30 neighbours form rows of the algebraic
matrices:
\begin{equation}
 \bm X,\bm W\in\mathbb R^{30\times27},\qquad
 \bm M=\bm X^{\mathsf T}\bm W\in\mathbb R^{27\times27}.
\label{eq:project-shapes}
\end{equation}
The centre is not a thirty-first row. After the 30 neighbour weights are obtained, its
weight is set to $w_0=1-\sum_{n=1}^{30}w_n$ so that constants are reproduced exactly.

\begin{center}\small
\begin{tabularx}{\textwidth}{@{}P{3.0cm}P{5.4cm}X@{}}
\toprule
\textbf{Where the selector stops} & \textbf{Points forced into the final footprint}
& \textbf{What may change} \\
\midrule
Compact tier & All 30 noncentre points of the fixed pyramid & Only the scale parameters
change; no point is replaced. \\
Repair source A & The same 30 pyramid neighbours & The amplification ceiling is relaxed. \\
Repair source B & 24--27 QRCP anchors from the pyramid & The remaining 6--3 points are
replaced by eligible non-template candidates. \\
Repair source C & 27 QRCP anchors from a pool prefix & Three additional eligible
non-anchor candidates complete the footprint. \\
\bottomrule
\end{tabularx}
\end{center}

Therefore the protected pool window, the nearest points in $\rho$, and the final
30-neighbour footprint are different sets. Across repaired nodes there is no universal
list of 30 index offsets.

\subsubsection{Compact tier: fixed index template, then physical-coordinate test}
\label{sec:project-compact-tier}

\paragraph{Purpose of this stage.}
This stage proposes one compact 31-point topology in structured-index space, maps it to
the actual physical coordinates, and accepts the first scale choice that passes the
algebraic certificate. The topology proposes; it does not prove solvability.

\paragraph{Step 1: construct the template in index space.}
The first trial ris not obtained by taking the 30 nearest points in physical distance.
Instead, \texttt{b2\_pattern} chooses seven consecutive physical $k$-index layers. Near a
wall the entire seven-layer block is shifted into the physical domain. The resulting
layer offsets $d_k$ are sorted by increasing $|d_k|$; at equal distance the negative
offset is placed first. The numbers of points assigned to the seven ordered layers are
\begin{equation}
 (7,7,5,5,3,3,1),
 \label{eq:pyramid-counts}
\end{equation}
where each layer uses the streamwise offset order
$d_j=0,+1,-1,+2,-2,\ldots$. The counts sum to 31 and include the centre
$(d_j,d_k)=(0,0)$ once; removing that centre leaves the required 30 neighbours.

For an interior centre, the seven $d_k$ values are
\[
 0,-1,+1,-2,+2,-3,+3.
\]
At the lower wall they are $0,+1,\ldots,+6$, and at the upper wall they are
$0,-1,\ldots,-6$. Thus the template is fixed in structured index space. After the index
offsets are mapped through $Y(j_p,k_p),Z(j_p,k_p)$, however, the physical cloud can be
stretched, sheared or curved; it need not look like a regular pyramid or diamond.

The seven-layer topology is intended to avoid the observed failure in which a
nearest-point cloud collapses onto at most six nearly parallel layers. Six such physical
layers can lie on the zero set of a nonzero degree-six product polynomial, making the
$P_6$ sampling matrix deficient. The template removes that specific geometry in the
intended grids, but it is not an analytic proof that every mapped template is unisolvent.

\paragraph{Step 2: define the anisotropic coordinate scaling.}
Using the local spacings in \eqref{eq:project-local-spacing}, define
\begin{equation}
 A=\operatorname{clamp}(s_y/s_z,1/4096,4096),\qquad
 \beta\in\{0,\tfrac12,1\},\qquad
 a_y=A^{\beta/2},\quad a_z=a_y^{-1},
\label{eq:beta-ladder}
\end{equation}
and, for neighbour $n$, define the scaled physical displacements
\begin{equation}
 \overline Y_n=\frac{\Delta Y_n}{a_y},\qquad
 \overline Z_n=\frac{\Delta Z_n}{a_z},\qquad
 \overline r_n=\sqrt{\overline Y_n^2+\overline Z_n^2}.
\label{eq:project-scaled-displacement}
\end{equation}
The bars denote coordinate scaling; they do not denote an average. Because
$a_ya_z=1$, this is an invertible diagonal change of coordinates and does not change the
polynomial space that is reproduced.

\paragraph{Step 3: define $h$ and the Wendland support.}
For each value of $\beta$, the template length scale is the root-mean-square radius of
its 30 neighbours:
\begin{equation}
 h_{\rm tmpl}=
 \left(\frac1{30}\sum_{n=1}^{30}\overline r_n^2\right)^{1/2}.
\label{eq:project-template-h}
\end{equation}
Hence $h$ is \emph{not} the minimum grid spacing. The code tries
\[
 h=m_h\,h_{\rm tmpl},\qquad m_h\in\{1,1.2,1.5,2\},
\]
in the written order. The dimensionless radial coordinate of neighbour $n$ is
\begin{equation}
 q_n=\frac{\overline r_n}{h}
 =\frac{\sqrt{(\Delta Y_n/a_y)^2+(\Delta Z_n/a_z)^2}}{h}.
\label{eq:active-support}
\end{equation}
The radial factor is the compactly supported Wendland $C^2$ function
\begin{equation}
 \psi(q)=
 \begin{cases}
  (1-q/2)^4(1+2q), & 0\le q<2,\\
  0, & q\ge2.
 \end{cases}
\label{eq:project-wendland-c2}
\end{equation}
All 30 template neighbours must satisfy $q_n<2$; otherwise at least one row of
$\bm W$ is zero and the trial is rejected.

\paragraph{Step 4: compact-first acceptance.}
The loop order is $\beta=0,\tfrac12,1$, with the four $h$ multipliers tested inside each
$\beta$. The first pair $(\beta,h)$ that passes every gate in
\S\ref{sec:project-admissibility} is accepted immediately. Therefore the compact tier
does not optimise over alternative point sets or search for the smallest condition
number. It preserves the index-space template whenever that template is numerically
acceptable.

\subsubsection{Admissibility: rank, conditioning and amplification}
\label{sec:project-admissibility}

\paragraph{Purpose of this stage.}
This stage answers the central solvability question. No argument based only on the
pyramid shape is used to declare $\bm M$ nonsingular. Instead, $\bm X$, $\bm W$ and
$\bm M$ are assembled from the 30 mapped physical points and certified numerically.

\paragraph{Step 1: define the 27 basis labels.}
Let
\begin{equation}
 \mathcal A_6=
 \{(a,b)\in\mathbb{N}_0^2:1\le a+b\le6\}.
\label{eq:project-basis-set}
\end{equation}
This set contains $2+3+4+5+6+7=27$ exponent pairs. The constant pair $(0,0)$ is omitted
because the centre weight later enforces constant reproduction. We use
$\mu=\mu(a,b)$ for a column label of $\bm X$ and
$\alpha=\alpha(u,v)$ for a column label of $\bm W$, where both $(a,b)$ and $(u,v)$
independently range over $\mathcal A_6$. The two labels each have 27 possible values,
but $\mu$ is not set equal to $\alpha$. More explicitly, $\mu(a,b)$ and
$\alpha(u,v)$ mean ``the integer column occupied by this exponent pair'' in the fixed
code ordering; they are column-address maps, not additional polynomial variables.

For neighbour $n$, divide the scaled displacements in
\eqref{eq:project-scaled-displacement} by $h$:
\begin{equation}
 \xi_n=\frac{\overline Y_n}{h}=\frac{\Delta Y_n}{a_yh},\qquad
 \eta_n=\frac{\overline Z_n}{h}=\frac{\Delta Z_n}{a_zh},\qquad
 q_n=\sqrt{\xi_n^2+\eta_n^2}.
\label{eq:project-normalised-coordinates}
\end{equation}
Here $q_n$ is the Wendland radial coordinate of a \emph{stencil point}. It is not an LBM
direction label; lattice directions will be denoted by $d$ in
\S\ref{sec:project-final-solve}.

\paragraph{Step 2: define every entry of $\bm X$, $\bm W$ and $\bm M$.}
For $\mu=\mu(a,b)$, the Taylor row is
\begin{equation}
 X_{n,\mu(a,b)}=\frac{\xi_n^a\eta_n^b}{a!\,b!}.
\label{eq:project-X-entry}
\end{equation}
Let $H_r$ be the physicists' Hermite polynomial,
$H_0(x)=1$, $H_1(x)=2x$, and
$H_{r+1}(x)=2xH_r(x)-2rH_{r-1}(x)$. For
$\alpha=\alpha(u,v)$, the Hermite--Wendland row is
\begin{equation}
 W_{n,\alpha(u,v)}
 =\psi(q_n)\,2^{-(u+v)/2}
 H_u\!\left(\frac{\xi_n}{\sqrt2}\right)
 H_v\!\left(\frac{\eta_n}{\sqrt2}\right),
\label{eq:project-W-entry}
\end{equation}
where $\psi$ is given by \eqref{eq:project-wendland-c2}. The moment matrix is
\begin{equation}
 M_{\mu\alpha}
 =\sum_{n=1}^{30}X_{n\mu}W_{n\alpha},
 \qquad \bm M=\bm X^{\mathsf T}\bm W.
\label{eq:cross-moment-entry}
\end{equation}
This is a sum, not an average; dividing it by 30 would only multiply the whole matrix by a
nonzero scalar and would not change its rank.

\paragraph{Why the template does not prove that $\bm M$ is full rank.}
Let $\bm x_n$ and $\bm w_n$ be row $n$ of $\bm X$ and $\bm W$. Then
\begin{equation}
 \bm M=\sum_{n=1}^{30}\bm x_n^{\mathsf T}\bm w_n.
\label{eq:rank-one-sum}
\end{equation}
Each point contributes a matrix of rank at most one. Thirty contributions do not
automatically span 27 independent directions, and different cross terms can cancel.
Moreover, $\bm w_n$ is not a positive scalar multiple of $\bm x_n$. For example,
\begin{equation}
 X_{n,\mu(2,0)}=\frac{\xi_n^2}{2},\qquad
 W_{n,\alpha(2,0)}=\psi(q_n)(\xi_n^2-1),\qquad
 W_{n,\alpha(3,0)}=\psi(q_n)(\xi_n^3-3\xi_n).
\label{eq:hermite-not-proportional}
\end{equation}
Therefore $\bm M$ is not a positive weighted Gram matrix
$\bm X^{\mathsf T}\bm D\bm X$, and positive definiteness cannot be inferred from
full column rank of $\bm X$.

An equivalent exact statement follows from thin QR factorizations
\begin{equation}
 \bm X=\bm Q_X\bm R_X,\qquad
 \bm W=\bm Q_W\bm R_W,\qquad
 \bm M=\bm R_X^{\mathsf T}
       (\bm Q_X^{\mathsf T}\bm Q_W)\bm R_W.
\label{eq:cross-subspace-factorisation}
\end{equation}
Here the columns of $\bm Q_X,\bm Q_W\in\mathbb R^{30\times27}$ are orthonormal. If
$\bm X$ and $\bm W$ have full column rank, the triangular factors $\bm R_X,\bm R_W$ are
nonsingular, and
\begin{equation}
 \operatorname{rank}(\bm M)=27
 \quad\Longleftrightarrow\quad
 \operatorname{rank}(\bm Q_X^{\mathsf T}\bm Q_W)=27
 \quad\Longleftrightarrow\quad
 \operatorname{col}(\bm W)\cap\ker(\bm X^{\mathsf T})=\{\bm0\}.
\label{eq:cross-pairing-condition}
\end{equation}
The fixed pyramid has no proven uniform lower bound for this cross-space pairing on every
allowed warped grid. It must therefore be checked after its actual coordinates are known.

\paragraph{Step 3: exact uniqueness condition.}
For a right-hand side $\bm C\in\mathbb R^{27}$, the coefficient vector
$\bm\Psi\in\mathbb R^{27}$ is defined by $\bm M\bm\Psi=\bm C$. Its exact
unique-solvability condition is
\begin{equation}
 \operatorname{rank}(\bm M)
 =\operatorname{rank}([\bm M\mid\bm C])=27.
\label{eq:unique-solvability}
\end{equation}
Because $\bm M$ is a $27\times27$ square matrix, $\operatorname{rank}(\bm M)=27$ alone
implies the augmented-rank equality for every $\bm C$.

%% 作者在文章裡寫 "deterministic rank-revealing QR" 是在描述目的和性質，你的 code 裡實際跑的是 QRCP 這個具體演算法來達成這個目的。%%

\paragraph{Step 4: numerical-rank gates.}
For each trial, the decisive question is whether the assembled moment matrix
$\bm M=\bm X^{\mathsf T}\bm W$ has numerical rank 27. The code does not apply its QR
rank test directly to the raw $\bm M$. It first performs six fixed row/column max-norm
equilibration sweeps and forms
\begin{equation}
 \widetilde{\bm M}=\bm D_r\bm M\bm D_c,
 \qquad
 \operatorname{rank}(\widetilde{\bm M})=\operatorname{rank}(\bm M),
\label{eq:project-equilibrated-rank}
\end{equation}
The subscripts $r$ and $c$ mean \emph{row} and \emph{column}; they do not denote a
radius. The two matrices are
\begin{equation}
 \bm D_r=\operatorname{diag}(d^{\rm row}_0,\ldots,d^{\rm row}_{26}),
 \qquad
 \bm D_c=\operatorname{diag}(d^{\rm col}_0,\ldots,d^{\rm col}_{26}).
\label{eq:project-equilibration-diagonals}
\end{equation}
The source does not store 729 entries for either diagonal matrix. The arrays
\texttt{Dr[27]} and \texttt{Dc[27]} store only
$d_i^{\rm row}$ and $d_i^{\rm col}$, respectively. They start at one. During each
equilibration sweep, the code first divides every row by its largest absolute entry and
accumulates that multiplier in \texttt{Dr}; it then does the same for every column and
accumulates the multiplier in \texttt{Dc}. In index notation, one sweep performs
\begin{equation}
\begin{aligned}
 s_i^{\rm row}&=\max_j|\widetilde M_{ij}|,&
 \widetilde M_{i:}&\leftarrow\widetilde M_{i:}/s_i^{\rm row},&
 d_i^{\rm row}&\leftarrow d_i^{\rm row}/s_i^{\rm row},\\
 s_j^{\rm col}&=\max_i|\widetilde M_{ij}|,&
 \widetilde M_{:j}&\leftarrow\widetilde M_{:j}/s_j^{\rm col},&
 d_j^{\rm col}&\leftarrow d_j^{\rm col}/s_j^{\rm col}.
\end{aligned}
\label{eq:project-equilibration-sweep}
\end{equation}
Thus $\bm D_r$ rescales the 27 equations (rows), whereas $\bm D_c$ rescales the 27
unknowns (columns). Both contain positive finite diagonal entries and are therefore
nonsingular, so the exact rank is unchanged.

The rank gate uses only $\widetilde{\bm M}$. When a right-hand side is later solved, the
same scaling is used consistently:
\begin{equation}
 \widetilde{\bm C}=\bm D_r\bm C,
 \qquad
 \widetilde{\bm M}\bm y=\widetilde{\bm C},
 \qquad
 \bm\Psi=\bm D_c\bm y.
\label{eq:project-equilibrated-system}
\end{equation}
Substitution gives
$\bm D_r\bm M\bm D_c\bm y=\bm D_r\bm C$, which is equivalent to
$\bm M\bm\Psi=\bm C$. Therefore equilibration changes only numerical scaling, not the
original solution. The reason for testing $\widetilde{\bm M}$ is that entries associated
with polynomial degrees one through six can have very different magnitudes; a QR test on
the raw matrix can mistake this scale imbalance for rank loss.

The code then applies deterministic rank-revealing QR to
$\widetilde{\bm M}^{\mathsf T}$. If $p_i$ is the residual norm of the $i$th QR pivot, it
is counted as an independent direction only when
\begin{equation}
 p_i>\tau_{\widetilde M}, \quad
  \tau_{\widetilde M}
 =27\,\epsilon_{\rm dbl}\,\|\widetilde{\bm M}\|_F.
\label{eq:project-rank-threshold}
\end{equation}
The result is therefore interpreted as follows:
\begin{itemize}[leftmargin=2.0em]
\item if QR finds 27 such pivots,
      $\operatorname{rank}_{\rm num}(\widetilde{\bm M})=27$ and the moment-matrix
      rank gate passes;
\item if QR finds fewer than 27, the current $(\text{stencil},\beta,h)$ trial is
      rejected and the selector proceeds to the next trial. No weights are exported.
\end{itemize}

Before this decisive test, the same QR routine also checks column-normalised versions of
$\bm X$ and $\bm W$. These two preliminary tests identify whether rank loss already
originates in polynomial sampling or in the Hermite--Wendland rows. Consequently every
trial must pass
\begin{equation}
 \operatorname{rank}_{\rm num}(\bm X)=27,\qquad
 \operatorname{rank}_{\rm num}(\bm W)=27,\qquad
 \operatorname{rank}_{\rm num}(\widetilde{\bm M})=27.
\label{eq:three-ranks}
\end{equation}
The first two conditions are necessary because
$\operatorname{rank}(\bm X^{\mathsf T}\bm W)
\le\min\{\operatorname{rank}(\bm X),\operatorname{rank}(\bm W)\}$, but they do not
guarantee the third. Hence the QR result for $\widetilde{\bm M}$ is the decisive rank
test. This is a double-precision numerical certificate, not a symbolic proof of
$\det(\bm M)\ne0$.

\paragraph{Step 5: reciprocal-conditioning gate.}
The equilibrated matrix must satisfy
\begin{equation}
 \operatorname{rcond}_{\infty}(\widetilde{\bm M})>\gamma_{27},\qquad
 \gamma_{27}=
 \frac{27\epsilon_{\rm dbl}}{1-27\epsilon_{\rm dbl}}
 \simeq5.9952\times10^{-15},
\label{eq:project-rcond}
\end{equation}
where $\epsilon_{\rm dbl}$ is double-precision machine epsilon. This is a
roundoff-level rejection threshold; it is not a claim that every accepted matrix has a
small condition number. Here
$\operatorname{rcond}_{\infty}(\bm B)=
1/(\|\bm B\|_\infty\|\bm B^{-1}\|_\infty)$ for a nonsingular $\bm B$.
The raw $\operatorname{rcond}_{\infty}(\bm M)$ is retained for diagnosis but is not an
acceptance gate.

\paragraph{Step 6: direction-independent amplification gate.}
Define two 27-component probe right-hand sides. The $Y$ probe
$\bm e^{[Y]}$ is one at basis label $(1,0)$ and zero elsewhere; the $Z$ probe
$\bm e^{[Z]}$ is one at $(0,1)$ and zero elsewhere. For
$\ell\in\{Y,Z\}$, solve
\begin{equation}
 \bm M\bm\Psi^{[\ell]}=\bm e^{[\ell]},\qquad
 w_n^{[\ell]}=\sum_{\alpha=1}^{27}
 W_{n\alpha}\Psi_\alpha^{[\ell]},\qquad
 w_0^{[\ell]}=1-\sum_{n=1}^{30}w_n^{[\ell]}.
\label{eq:project-probe-weights}
\end{equation}
The direction-independent amplification is
\begin{equation}
 \lambda_0=
 \max_{\ell\in\{Y,Z\}}\sum_{n=0}^{30}|w_n^{[\ell]}|.
\label{eq:lambda-zero}
\end{equation}
The compact tier requires $\lambda_0\le512$; repair trials require
$\lambda_0\le4096$. These limits reject extreme cancellation. They are ceilings, not
quantities minimised by the selector.

Passing \eqref{eq:three-ranks}--\eqref{eq:lambda-zero} is an a posteriori
finite-precision certificate for the particular coordinates, $h$ and anisotropic scales
that were tested. It is not an analytic proof that every stencil proposed by the
index-space pyramid is nonsingular. A failed compact trial is sent to the next
$(\beta,h)$ pair or to the repair tier; if every trial fails, no stencil is exported.

\subsubsection{Repair tier and the leading-error objective}
\label{sec:project-repair-tier}

\paragraph{Purpose of this stage.}
This stage is a fallback, not a second optimisation pass over an already accepted
template. It constructs alternative 30-neighbour clouds only after all compact trials
fail, applies exactly the same solvability gates, and then ranks only the survivors by
their leading degree-seven moment leakage.

\paragraph{When repair is entered.}
Repair is not run after a compact stencil has already passed. It is entered only when
every combination
\[
 \beta\in\{0,\tfrac12,1\},\qquad
 h/h_{\rm tmpl}\in\{1,1.2,1.5,2\}
\]
for the fixed pyramid has failed at least one admissibility gate in
\S\ref{sec:project-admissibility}.

\paragraph{How the three repair sources construct 30 neighbours.}
Here QRCP means the deterministic rank-revealing QR procedure that treats physical points
as row candidates and orders the selected rows by their independent contribution to the
Taylor matrix. Let $N_{\rm pool}$ denote the number of retained pool points. In the
table, a prefix of nominal length $L$ means the first $\min(L,N_{\rm pool})$ points.
\begin{center}\small
\begin{tabularx}{\textwidth}{@{}P{2.2cm}P{5.5cm}X@{}}
\toprule
\textbf{Source} & \textbf{Construction} & \textbf{Points guaranteed in that trial} \\
\midrule
A: unchanged pyramid &
Use exactly the same 30 neighbours as the compact template, but apply the repair-tier
limit $\lambda_0\le4096$. &
All 30 template neighbours. \\
B: minimally modified pyramid &
Run QRCP on the 30 template rows. If its returned rank
$r_{\rm tmpl}$ is at least 24, retain the first
$\min(r_{\rm tmpl},27)$ template anchors. Exclude every other template point from refill,
then add eligible non-template candidates in pool order until 30 points are present. &
Between 24 and 27 template anchors; consequently 6 to 3 points are replaced. \\
C: pool-prefix cloud &
Open the first 40, 60, 100 and 160 pool entries in turn. Among entries satisfying
$q_n<2$, QRCP must find exactly 27 independent anchors; the first three eligible
non-anchors in pool order complete the cloud. &
Exactly 27 pool-QRCP anchors and three deterministic local fill points. \\
\bottomrule
\end{tabularx}
\end{center}
Thus source B can replace three to six points. The nearby source-code comment that says
at most three replacements is stale relative to the executable condition
$r_{\rm tmpl}\ge24$.

\paragraph{Scale parameters tried by repair.}
Sources A and B reuse $h_{\rm tmpl}$ from \eqref{eq:project-template-h} and try
\[
 h/h_{\rm tmpl}\in\{1,1.2,1.5,2\}.
\]
For source C, let $h_{\rm base}$ be the RMS value in
\eqref{eq:project-template-h} computed instead from the first 30 candidate-pool entries
at the current $\beta$. Source C tries
\[
 h/h_{\rm base}\in\{1.2,1.5,2,2.5,3,4,5,6\}.
\]
All three sources are generated for each $\beta$ in \eqref{eq:beta-ladder}, and every
complete trial must independently pass the support, rank, reciprocal-conditioning and
$\lambda_0$ gates of \S\ref{sec:project-admissibility}.

\paragraph{How the winning repair cloud is chosen.}
The first-moment probe weights $w_n^{[Y]}$ and $w_n^{[Z]}$ were defined in
\eqref{eq:project-probe-weights}. Using the scaled physical displacements
$\overline Y_n,\overline Z_n$ from \eqref{eq:project-scaled-displacement}, define
\begin{equation}
 T_7=
 \sum_{\ell\in\{Y,Z\}}\ \sum_{\substack{a+b=7\\a,b\in\mathbb N_0}}
 \frac{1}{a!\,b!}
 \left|
 \sum_{n=1}^{30}
 w_n^{[\ell]}\overline Y_n^a\overline Z_n^b
 \right|.
\label{eq:T7-selector}
\end{equation}
The $P_6$ moment conditions control degrees zero through six, so degree seven is the first
uncontrolled Taylor degree. The code evaluates \eqref{eq:T7-selector} for every
admissible source-A, source-B and source-C trial and retains the strictly smallest value.
Because the comparison is strict, an exact tie retains the trial encountered first in the
fixed loop order.

$T_7$ is a direction-independent comparison proxy based on the two first-moment
responses. It is not the complete truncation error for a particular LBM departure point,
and it is meaningful here only for comparing trials at the same arrival node. The older
helper proportional to $\lambda_0r_{\max}^7$ has no caller in the active selector and is
not the current optimisation objective.

\subsubsection{Final direction-dependent solve and fail-closed export}
\label{sec:project-final-solve}

\paragraph{Purpose of this stage.}
The preceding stages choose only the geometry. This final stage supplies one departure
point at a time, computes that direction's weights on the already selected footprint,
and prevents any failed solution from reaching the time-step loop.

\paragraph{One footprint, different weights for different lattice directions.}
Stages 1--4 select geometry without using an LBM velocity. Consequently one arrival node
has one set of 31 storage IDs shared by all lattice directions. Direction-specific data
enter only now through the departure point.

Let $d$ denote an LBM direction and let
$(e_y^{(d)},e_z^{(d)})$ be its two lattice-velocity components. The physical displacement
from the arrival node to the departure point is, with $\Delta t$ denoting the global
time step,
\begin{equation}
 \delta_y^{(d)}=-\Delta t\,e_y^{(d)},\qquad
 \delta_z^{(d)}=-\Delta t\,e_z^{(d)}.
\label{eq:project-departure-displacement}
\end{equation}
Use the same $a_y,a_z,h$ as the accepted stencil and define
\begin{equation}
 \xi_\delta^{(d)}=\frac{\delta_y^{(d)}}{a_yh},\qquad
 \eta_\delta^{(d)}=\frac{\delta_z^{(d)}}{a_zh}.
\label{eq:project-departure-normalised}
\end{equation}
For the Taylor label $\mu=\mu(a,b)$ introduced in
\eqref{eq:project-basis-set}, the right-hand side is
\begin{equation}
 C_{\mu(a,b)}^{(d)}
 =\frac{(\xi_\delta^{(d)})^a(\eta_\delta^{(d)})^b}{a!\,b!}.
\label{eq:project-direction-rhs}
\end{equation}
This is the precise meaning of the abbreviated statement
$\bm C^{(d)}=\bm X(\bm\delta^{(d)})$.

\paragraph{Solve for coefficients, then recover interpolation weights.}
For direction $d$, solve
\begin{equation}
 \bm M\bm\Psi^{(d)}=\bm C^{(d)}.
\label{eq:project-direction-solve}
\end{equation}
The 30 neighbour weights and the structural centre weight are then
\begin{equation}
 w_n^{(d)}=\sum_{\alpha=1}^{27}
 W_{n\alpha}\Psi_\alpha^{(d)}
 \quad (n=1,\ldots,30),\qquad
 w_0^{(d)}=1-\sum_{n=1}^{30}w_n^{(d)}.
\label{eq:project-direction-weights}
\end{equation}
Thus the footprint and $\bm M$ are shared across directions, whereas
$\bm C^{(d)}$, $\bm\Psi^{(d)}$ and $\bm w^{(d)}$ depend on the departure direction.

\paragraph{Hard solve gates.}
The code equilibrates \eqref{eq:project-direction-solve} and solves it independently by
GEPP, explicit PLU and Gauss--Jordan. All three solutions undergo double-double iterative
refinement using the factored residual of the original product
$\bm X^{\mathsf T}\bm W$. Export is permitted only if
\begin{enumerate}[label=(\roman*),leftmargin=2.0em]
\item all three solvers report success;
\item each solver's normwise and componentwise backward errors are at most
      $64\gamma_{27}$;
\item $\operatorname{rcond}_{\infty}(\widetilde{\bm M})>\gamma_{27}$; and
\item for every component of $\bm\Psi$, each solver pair $(s,t)$ satisfies
      $|\Psi_i^{(s)}-\Psi_i^{(t)}|/
       \max(1,|\Psi_i^{(s)}|,|\Psi_i^{(t)}|)\le64\gamma_{27}$.
\end{enumerate}
In item (iv), $i=1,\ldots,27$, and $(s,t)$ runs over the three pairs
GEPP--PLU, GEPP--Gauss--Jordan and PLU--Gauss--Jordan.
Only after these tests does the code form \eqref{eq:project-direction-weights} and allow
\texttt{labfm\_compact\_export} to copy the 31 IDs and weights into the time-step data.

\paragraph{Warnings versus hard failures.}
After a successful solve, let
$\bm w_{\rm nb}^{(d)}=(w_1^{(d)},\ldots,w_{30}^{(d)})^{\mathsf T}$ denote the neighbour
weights without the centre weight. The code also reports
\begin{equation}
 R_d=
 \frac{\|\bm X^{\mathsf T}\bm w_{\rm nb}^{(d)}-\bm C^{(d)}\|_2}
      {\|\bm C^{(d)}\|_2},
 \qquad
 \Lambda_d=\sum_{n=0}^{30}|w_n^{(d)}|,
\label{eq:project-direction-diagnostics}
\end{equation}
with the zero denominator treated as one for the rest direction. The thresholds
$R_d>10^{-8}$ and $\Lambda_d>6$ currently set warning bits only; they are not the same as
the selection hard gates \eqref{eq:project-rcond} and \eqref{eq:lambda-zero}.

If point selection or any hard solve gate fails, no compact stencil is exported. The
active path has no pseudoinverse, regularisation, wall-ghost substitution or
nearest-neighbour fallback. A concise Chinese source map is maintained separately in
\path{labfm/stencil.md}.

\subsection{A parity diagnostic for sampling adequacy}
\label{sec:parity}

The angular sampling condition of \S\ref{sec:sampling} is stated as a geometric rule
and measured in~\cite{KLN2020} through the condition number, which is a lagging
indicator: by the time $\kappa$ jumps, undersampling is already severe. A direct and
inexpensive leading indicator follows from the parity structure of the continuum limit
\New.

For the Hermite construction, $\mathrm{He}_n$ has parity $(-1)^n$, so the integrand of
$\langle X^{\alpha},W^{\beta}\rangle$ with $\alpha=(a,b,c)$, $\beta=(a',b',c')$ is odd
in at least one coordinate unless every one of $a{+}a'$, $b{+}b'$, $c{+}c'$ is even. On
an isotropic, symmetric support the integral then vanishes:
\begin{equation}
\bigl\langle X^{\alpha},W^{\beta}\bigr\rangle=0
\quad\text{unless } a{+}a',\ b{+}b',\ c{+}c'\ \text{are all even.}
\label{eq:parity}
\end{equation}
This is the analytic content of the pattern of machine-zero entries displayed
in~\cite[Fig.~5]{KLN2020}. Define the \emph{parity defect}
\begin{equation}
\eta_i\;=\;\Biggl(
\frac{\sum_{(\alpha,\beta)\ \text{forbidden}}\bigl|M^{\alpha\beta}_i\bigr|^{2}}
{\sum_{\alpha,\beta}\bigl|M^{\alpha\beta}_i\bigr|^{2}}\Biggr)^{1/2},
\label{eq:parity-defect}
\end{equation}
computable in $O(p^{2})$ with no quadrature. It serves three purposes: on a uniform
Cartesian cloud $\eta_i\lesssim10^{-14}$ is a sharp test of the assembly (a failure
indicates a multi-index ordering or Hermite recurrence error); the value of $h/s$ at
which $\eta_i$ begins to rise locates the critical ratio ahead of any change in
$\kappa$; and a large $\eta_i$ at a node correctly flags an anisotropic or one-sided
stencil, identifying where the order must be reduced.

Where an absolute rather than relative measure is wanted, the continuum reference
$N_i\langle X^{\alpha},W^{\beta}\rangle$ of \eqref{eq:sampling-norm} may be evaluated
by quadrature over the support and compared with the assembled $\bm M_i$.

\subsection{Published stability results}
\label{sec:stab-pub}

For completeness, and because they bound what may legitimately be assumed:
\begin{itemize}
\item On ideal Cartesian distributions the convective operator has
$\Re(\lambda)\approx0$ provided $\mathrm dW^{0}/\mathrm dr=0$ on the support boundary,
a condition satisfied by the quadratic and Wendland fundamental functions but not by
the conic one~\cite[\S 4.2]{KLN2020}. On distributions carrying disorder,
$\Re(\lambda)>0$ for all choices tested.
\item The Laplacian operator is stable, but the admissible level of distribution noise
decreases as $m$ increases~\cite[\S 4.2, Table 1]{KLN2020}.
\item Hyperviscosity operators up to $\nabla^{m}$ may be generated from the same
system \eqref{eq:moments} by an appropriate choice of $\bm C^{\rm hv}$
\cite[Eq.~(25)]{KLN2020}.
\end{itemize}

\begin{remark}
These results concern \emph{differential} operators and therefore do not transfer
directly to the present scheme, in which streaming is effected by interpolation. The
governing stability condition for the scheme specified here is
\eqref{eq:vn-crit} below.
\end{remark}

\subsection{Boundaries: what the literature actually reports}
\label{sec:boundary}

This subsection supersedes the abstract of~\cite{KLN2020} and is the single most
consequential correction in this revision.

\paragraph{The 2020 claim.} \cite{KLN2020} states in its abstract that ``at domain
boundaries (with incomplete support) LABFM automatically provides one-sided differences
of the same order as the internal scheme, up to 4th order'', and in \cite[\S
IV\,A]{KLN2020} that convergence at order $m$ is retained near boundaries provided the
distribution noise is small enough and $h/s$ large enough.

\paragraph{The 2021 retraction.} \cite[\S 7.1]{KL2021}, by the same authors, reports
that this does not hold:
\begin{itemize}
\item for one-sided stencils and $m>4$, ``the conditioning of the linear systems $M_i$
was such that the solver didn't yield a valid solution'';
\item ``a new stability analysis since \cite{KLN2020} (via the eigenvalues of the global
derivative operator) has shown that for one-sided LABFM with $m=4$ the eigenvalues
spread along the positive real axis: \emph{the operators are unstable}'';
\item for nearly one-sided stencils one node in from the boundary, the viscous operators
are stable but the convective operators remain unstable.
\end{itemize}
\cite{KL2021} accordingly \emph{abandons} one-sided LABFM and adopts a hybrid, which
requires a strip of five uniformly distributed, locally orthogonal nodes along every
boundary:

\begin{center}
\begin{tabular}{@{}llll@{}}
\toprule
\textbf{Node} & $\partial/\partial\xi$ \textbf{(normal)}
& $\partial/\partial\eta$ \textbf{(tangential)} & \textbf{Second derivatives} \\
\midrule
$i_0$ (on boundary) & one-sided FD & 1D-LABFM & FD $+$ 1D-LABFM \\
$i_1$ & asymmetric FD & 1D-LABFM & LABFM, $m=4$ \\
$i_2$ & centred FD & 1D-LABFM & LABFM, $m=4$ \\
$i_3$, $i_4$, interior & LABFM & LABFM & LABFM \\
\bottomrule
\end{tabular}
\end{center}
\noindent(\cite[Table 1]{KL2021}; the one-dimensional formulation used for tangential
derivatives is given in \cite[Appendix A, Eqs.~(A1)--(A7)]{KL2021}, and first
derivatives at $i_1,i_2$ are rotated back to Cartesian components through the Jacobian
\cite[Eq.~(24)]{KL2021}.) The production configuration of~\cite{KL2021} is therefore
$m=6$ in the interior and an effectively fourth-order treatment in the boundary strip.

\paragraph{Consequences for the present scheme.} Four, and they must be kept distinct.
\begin{enumerate}[label=(\roman*)]
\item \textbf{The conditioning risk transfers, but the categorical stencil conclusion
does not.} For a fixed cloud, basis and scaling, $\bm M_i$ is independent of the
right-hand side, so a singular derivative construction cannot be rescued merely by
asking it to interpolate. This does not prove that every other degree-six physical
cloud is singular. The current selector of \S\ref{sec:project-stencil} changes the cloud,
permits a controlled anisotropic coordinate scaling, and requires the actual one-sided
$\bm X$, $\bm W$ and equilibrated $\bm M$ to pass rank and reciprocal-condition gates.
It therefore treats the 2021 result as a fail-closed warning to test each wall cloud,
not as a theorem that $m>4$ is impossible for all interpolation selectors.
\item \textbf{The instability result does not transfer directly.} It is an eigenvalue
statement about the global \emph{derivative} matrix. The present scheme streams by
interpolation, whose stability is governed by \eqref{eq:vn-crit}. Whether one-sided
interpolation operators are admissible must therefore be settled by measuring
$\max_{\bm k}|G_q(\bm k)|$ on one-sided stencils, not inferred from~\cite{KL2021}.
\item \textbf{The current LABFM footprint is physical and one-sided at a wall.}
Wall-normal ghosts are forbidden by \eqref{eq:enum-set}; the seven-layer template is
clamped into the domain. Ghost layers remain necessary for the incumbent ITB/GILBM
paths, boundary closure and statistics contract, but they are not LABFM stencil members.
If the finite-difference hybrid of~\cite{KL2021} is later adopted, it is a different
method branch rather than the selector presently encoded in
\path{Lagrange_vs_LABFM/2.labfm/}.
\item \textbf{Local algebraic acceptance is not global stability or physical
validation.} Passing \eqref{eq:three-ranks}--\eqref{eq:lambda-zero} and the three-solver
gate establishes that the local interpolation weights are numerically constructible.
It does not establish the one-sided global amplification bound, observed seventh-order
convergence at every wall geometry, or correctness of a wall-resolved production run.
\end{enumerate}
Open item U4 records what must be measured before a wall-resolved production claim is
made.

%========================================================================
\section{The departure-point interpolation operator}
\label{sec:interp}

\subsection{Definition}

Because the characteristics of the free-streaming operator are straight lines in
physical space, the departure point is available in closed form,
\begin{equation}
\bm\delta_q=-\,\bm e_q\,c\,\Delta t,
\qquad
\bm r^{*}_q=\bm r_i+\bm\delta_q ,
\label{eq:departure}
\end{equation}
with $c$ the lattice speed and $\Delta t$ the global time step. No integration of the
characteristic and no metric information is required.

\begin{remark}[The three points, and why the departure point carries no index]
\label{rem:threepoints}
Three locations appear and are easily conflated. The node $\bm r_i$ is the Taylor
origin, the stencil centre, and the arrival node at which the result is stored; its
role is exactly that of $i$ in~\cite{KLN2020} and is unchanged. The nodes $\bm r_j$,
$j\in\mathcal N_i$, supply data; they are neighbours, and \emph{none of them is the
departure point}. The departure point $\bm r^{*}_q$ is neither: it is a query location
that enters only through the right-hand side of \eqref{eq:interp-op}. It does not
appear in $\bm M_i$ at all, which is precisely why one factorisation serves all
nineteen discrete velocities (Remark~\ref{rem:shared}).
\end{remark}

\begin{proposition}[Admissibility of the node-centred stencil; \New]
\label{prop:cfl}
Centring the stencil at $\bm r_i$ requires the departure point to lie within the
support, $|\bm\delta_q|\le2h_i$. With $\Delta t=\mathrm{CFL}\cdot s_{\min}/c$ and
$\max_q|\bm e_q|=\sqrt2$ for D3Q19, and taking the worst case $s_i=s_{\min}$,
\begin{equation}
\mathrm{CFL}\;\le\;\sqrt2\,\bigl(h/s\bigr).
\label{eq:cfl-bound}
\end{equation}
\end{proposition}
\noindent At $h/s=1.6$ this gives $\mathrm{CFL}\le2.26$; a working margin of one half of
the support radius halves it. The present operating point, $\mathrm{CFL}=0.5$, places
the departure point at between $2\,\%$ and $22\,\%$ of the support radius depending on
local spacing, comfortably inside. The bound is recorded because it is not independent
of the time-step question: the relaxation-time constraint
$\tau=3\nu/\Delta t+\tfrac12$ becomes binding at a comparable value, so any attempt to
enlarge $\Delta t$ will approach both limits at once. Beyond \eqref{eq:cfl-bound} the
stencil would have to be centred on the departure point, at the cost of one moment
matrix per discrete velocity (Remark~\ref{rem:origins}).

\begin{definition}[Departure-point interpolation operator; \New]
\label{def:interp}
Using the same moment matrix $\bm M_i$ as in \eqref{eq:moments}, define for each
discrete velocity $q$
\begin{equation}
\bm M_i\,\bm\Psi^{q}_i=\widetilde{\bm X}(\bm\delta_q)
\equiv\bm X(\bm\delta_q)\!\cdot\!\bm H_i ,
\qquad
w^{q}_{ji}=\bm W_{ji}\cdot\bm\Psi^{q}_i ,
\label{eq:interp-op}
\end{equation}
and take as the streaming step
\begin{equation}
f_q(\bm r_i,t+\Delta t)
=f^{\rm post}_{q,i}
+\sum_{j\in\mathcal N_i}\bigl(f^{\rm post}_{q,j}-f^{\rm post}_{q,i}\bigr)w^{q}_{ji}
=w^{q}_{\rm cen}f^{\rm post}_{q,i}+\sum_{j\in\mathcal N_i}w^{q}_{ji}f^{\rm post}_{q,j},
\label{eq:stream}
\end{equation}
where $w^{q}_{\rm cen}\equiv1-\sum_{j}w^{q}_{ji}$.
\end{definition}

The form of the right-hand side follows from applying \eqref{eq:taylor} to the
departure point, $\phi(\bm r^{*}_q)-\phi_i=\bm X(\bm\delta_q)\cdot\bm D\phi|_i$, and
comparing with \eqref{eq:Lop}. Consistency of the preconditioning is guaranteed by
using the same $\bm H_i$ in $\widetilde{\bm X}(\bm\delta_q)$ as in
$\widetilde{\bm X}_{ji}$.

\subsection{Polynomial reproduction}

\begin{theorem}[Exact reproduction at the departure point; \New]
\label{thm:exact}
Let $\bm M_i$ be non-singular, let $\bm X$ contain all monomials of total degree at
most $m$, and let $\bm\Psi^{q}_i$ solve \eqref{eq:interp-op}. Then for every polynomial
$P$ of total degree at most $m$,
\begin{equation*}
P_i+\sum_{j\in\mathcal N_i}(P_j-P_i)\,w^{q}_{ji}=P(\bm r_i+\bm\delta_q),
\end{equation*}
exactly, with no truncation error.
\end{theorem}

\begin{proof}
For a polynomial of total degree at most $m$ the truncated expansion \eqref{eq:taylor}
is an identity rather than an approximation, so that $P_j-P_i=\bm X_{ji}\cdot\bm DP|_i$.
Hence
\begin{align*}
\sum_{j}(P_j-P_i)w^{q}_{ji}
&=\sum_{j}\bigl(\bm X_{ji}\cdot\bm DP|_i\bigr)\bigl(\bm W_{ji}\cdot\bm\Psi^{q}_i\bigr)
=\bm DP|_i\cdot\Bigl(\sum_{j}\bm X_{ji}\otimes\bm W_{ji}\Bigr)\bm\Psi^{q}_i\\
&=\bm DP|_i\cdot\bm M_i\bm\Psi^{q}_i
=\bm DP|_i\cdot\bm X(\bm\delta_q)
=P(\bm r_i+\bm\delta_q)-P(\bm r_i),
\end{align*}
the last step again by the exactness of \eqref{eq:taylor} for polynomials of degree at
most $m$. Adding $P_i$ gives the result. The preconditioned form is equivalent, since
both sides of the linear system are multiplied by the same diagonal matrix.
\end{proof}

\begin{corollary}[Exact preservation of constants]
\label{cor:const}
Taking $P\equiv\text{const}$ gives $P_j-P_i\equiv0$, so that \eqref{eq:stream} returns
$f_{q,i}$ identically. A uniform field therefore generates no spurious mass source.
This follows from the difference form $(\cdot)_{ji}$ alone and requires no additional
partition-of-unity constraint on the weights.
\end{corollary}

\begin{corollary}[Degenerate zero-displacement limit]
\label{cor:zero}
If $\bm\delta_q=\bm 0$ then $\bm X(\bm 0)=\bm 0$, hence $\bm\Psi^{q}_i=\bm 0$ and
$w^{q}_{ji}\equiv0$, and \eqref{eq:stream} reduces to the identity exactly. This
provides the sharpest available unit test of the assembled operator (rung L5a of
\S\ref{sec:ladder}).
\end{corollary}

\begin{remark}[The operator is approximating, not interpolating]
\label{rem:noninterp}
Lagrange interpolation reproduces nodal values exactly: a departure point coinciding
with a node returns that node's value. The operator of Definition~\ref{def:interp} does
not have this property, since it enforces polynomial consistency rather than nodal
collocation, in common with moving least squares approximation. Consequently the
condition ``the scheme must recover the standard lattice Boltzmann method when the
departure point coincides with a node'' is \emph{not} a valid verification criterion
for this scheme. The only exact degeneracy is Corollary~\ref{cor:zero}.
\end{remark}

\subsection{Truncation error and the amplification constant}

For a general smooth field, $f_j-f_i=\bm X_{ji}\cdot\bm Df|_i+R_{ji}$ with
$R_{ji}=O(|\rji|^{m+1})$. Theorem~\ref{thm:exact} then gives
\begin{equation}
\varepsilon^{q}_i
=\Bigl|\sum_{j}R_{ji}w^{q}_{ji}-R(\bm\delta_q)\Bigr|
\le\Lambda^{q}_i\max_{j}|R_{ji}|+|R(\bm\delta_q)|
=O\!\bigl(\Lambda^{q}_i\,h_i^{\,m+1}\bigr),
\label{eq:trunc}
\end{equation}
where
\begin{equation}
\boxed{\ \Lambda^{q}_i\equiv\sum_{j\in\mathcal N_i}\bigl|w^{q}_{ji}\bigr|\ }
\label{eq:lebesgue}
\end{equation}
plays the role of a Lebesgue constant for the operator \New. It bounds both the error
amplification in \eqref{eq:trunc} and the growth of the streaming step in the maximum
norm, since \eqref{eq:stream} gives
$|f^{\rm new}|\le\bigl(|w^{q}_{\rm cen}|+\Lambda^{q}_i\bigr)\max_j|f_j|$. The quantity
$\Lambda^{q}_i$ is to be computed and reported for every node and every discrete
velocity during preprocessing; it is the second diagnostic, alongside the condition
number, on which acceptance of a stencil is conditioned.

\subsection{Von Neumann amplification and the stability criterion}
\label{sec:vn}

Substituting a plane wave $f=\exp(\mathrm i\bm k\cdot\bm r)$ into \eqref{eq:stream}
gives the amplification factor
\begin{equation}
G_q(\bm k)=1+\sum_{j\in\mathcal N_i}
\Bigl(e^{\mathrm i\bm k\cdot\rji}-1\Bigr)w^{q}_{ji},
\label{eq:amp}
\end{equation}
whose exact counterpart is the pure phase shift $e^{\mathrm i\bm k\cdot\bm\delta_q}$,
of unit modulus. Hence:

\begin{center}
\fbox{\parbox{0.95\textwidth}{
\New\ \textbf{Stability criterion for the streaming step.}
\begin{equation}
\max_{|\bm k|\le k_{\rm Ny}}\bigl|G_q(\bm k)\bigr|\le1+\epsilon_{\rm tol},
\qquad q=0,\dots,18,
\qquad k_{\rm Ny}=\pi/s_i .
\label{eq:vn-crit}
\end{equation}
The criterion is evaluated per node during preprocessing. A node violating it must
have its polynomial order reduced or its stencil enlarged. The phase error
$\arg G_q(\bm k)-\bm k\cdot\bm\delta_q$ is recorded separately as the dispersion
diagnostic.
}}
\end{center}

\begin{remark}[Relation to the source]
Criterion \eqref{eq:vn-crit} is an adaptation of steps 3 and 4 of the stencil
optimisation algorithm of~\cite[\S 4]{KL2021}, in which the amplitude response of the
\emph{differential} operators is evaluated at $n_k=16$ wavenumbers up to the Nyquist
limit \cite[Eqs.~(13a)--(13c)]{KL2021} and the stencil is accepted only while
$\max(A^{x},A^{y},A^{L})\le1.01$. The present criterion replaces the effective
wavenumber of a derivative by the amplification factor of a translation; the tolerance
$\epsilon_{\rm tol}=0.01$ and the sampling $n_k=16$ are retained from the source. The
same Fourier machinery underlies the resolving-power analysis
of~\cite{Broadley2025,Broadley2026}, which provides the natural route to a future
implicit (compact) variant.
\end{remark}

\subsection{Local affine rescaling for anisotropic node distributions}
\label{sec:affine}

Near a wall the node spacing normal to the surface is much smaller than the tangential
spacing, so the stencil is a flattened ellipsoid in physical space and $\bm M_i$ is
poorly conditioned. Let
$\bm S_i=\operatorname{diag}\bigl(h_{x,i}^{-1},h_{y,i}^{-1},h_{z,i}^{-1}\bigr)$ and
assemble and solve the system in the rescaled coordinates
$\tilde{\bm r}_{ji}=\bm S_i\rji$.

\begin{theorem}[Affine invariance of polynomial consistency; \New]
\label{thm:affine}
The space $\mathbb P_m$ of polynomials of total degree at most $m$ is invariant under
any invertible linear map. Consequently, enforcing degree-$m$ consistency in
$\tilde{\bm r}$ is equivalent to enforcing degree-$m$ consistency in $\bm r$.
\end{theorem}

\begin{proof}
Let $A$ be invertible and $P\in\mathbb P_m$. Then $P\circ A$ is a polynomial with
$\deg(P\circ A)=\deg P$, since each monomial of total degree $n$ maps to a linear
combination of monomials of total degree $n$. Theorem~\ref{thm:exact} applied to the
point set $\{A\rji\}$ with right-hand side $\bm X(A\bm\delta_q)$ is therefore
equivalent to the same statement for $\{\rji\}$ and $\bm X(\bm\delta_q)$.
\end{proof}

\begin{remark}[Implementation consequences]
Interpolation weights are dimensionless, so no inverse scaling of the weights is
required; it suffices to evaluate the right-hand side as
$\widetilde{\bm X}(\bm S_i\bm\delta_q)$. Weights of differential operators do carry
dimensions and must be rescaled by the corresponding powers of $\bm S_i$. The Hermite
ABFs \eqref{eq:hermiteABF} extend to three dimensions as
\begin{equation}
W^{\alpha}_{ji}=\frac{\psi(|\tilde{\bm r}_{ji}|)}{\sqrt{2^{\,a+b+c}}}\,
H_a\!\left(\frac{\tilde x_{ji}}{\sqrt2}\right)
H_b\!\left(\frac{\tilde y_{ji}}{\sqrt2}\right)
H_c\!\left(\frac{\tilde z_{ji}}{\sqrt2}\right),
\label{eq:hermite3d}
\end{equation}
for $\bm X$-entries with multi-index $\alpha=(a,b,c)$; this is the direct dimensional
extension of \cite[Eq.~(10)]{KL2021} \New. Throughout, a Greek superscript indexes the
\emph{basis}, running over $1,\dots,p$, while the Latin $q$ reserved in
\S\ref{sec:interp} indexes the \emph{discrete velocity}; the two are unrelated.
\end{remark}

%========================================================================
\section{A generalised moving least squares control formulation}
\label{sec:mls}

The construction of \S\ref{sec:interp} uses the LABFM machinery, but nothing in
Theorems~\ref{thm:exact} and~\ref{thm:affine}, in Corollaries~\ref{cor:const}
and~\ref{cor:zero}, or in the diagnostics \eqref{eq:lebesgue} and \eqref{eq:vn-crit}
requires it: they apply to any weight construction that enforces polynomial
consistency on a node cloud. The most established such alternative is moving least
squares~\cite{Lancaster1981} in its generalised, node-centred
form~\cite{Mirzaei2012,Trask2016}. This section specifies that alternative in a form
that is directly comparable, term by term, so that the selection between the two can
be made by measurement rather than by assertion.

\subsection{Definition}

Let $\bm q(\bm z)=[1,\,z_x,\,z_y,\,z_z,\,z_x^2/2,\,\dots]^{\mathsf T}$ be the complete
polynomial basis of total degree at most $m$, of dimension
$P=\binom{m+3}{3}=p_{\rm 3D}+1$, and let $\theta_{ji}=\theta(|\rji|/h_i)$ be a
compactly supported kernel. Write $\bm q_{ji}=\bm q(\bm S_i\rji)$ with $\bm S_i$ the
local affine scaling of \S\ref{sec:affine}. Define the Gram matrix and the shape
functions evaluated at the departure point,
\begin{equation}
\bm A_i=\sum_{j\in\mathcal N_i}\theta_{ji}\,\bm q_{ji}\bm q_{ji}^{\mathsf T},
\qquad
\Phi^{q}_{ji}=\bm q(\bm S_i\bm\delta_q)^{\mathsf T}\bm A_i^{-1}\theta_{ji}\bm q_{ji}.
\label{eq:mls}
\end{equation}

\begin{remark}[Why the kernel is centred on the node]
In classical moving least squares the kernel is centred on the \emph{evaluation}
point, so that $\bm A$ must be reassembled and refactorised for every target. With
nineteen targets per node this would forfeit the structural advantage identified in
Remark~\ref{rem:shared}. Centring the kernel on node $i$ instead --- the generalised
or fixed-kernel variant~\cite{Mirzaei2012,Trask2016} --- makes $\bm A_i$
target-independent, so that a single factorisation serves all nineteen departure
points, exactly as $\bm M_i$ does. Only this variant constitutes a fair control.
\end{remark}

Because $1\in\operatorname{span}\bm q$, the shape functions form a partition of unity,
$\sum_j\Phi^{q}_{ji}=1$, and the streaming step may be written in the same difference
form as \eqref{eq:stream} with
\begin{equation}
w^{q,\mathrm{MLS}}_{ji}=\Phi^{q}_{ji}.
\label{eq:mls-weights}
\end{equation}
The diagnostics \eqref{eq:lebesgue} and \eqref{eq:amp} therefore apply verbatim to
both arms, with identical definitions.

\subsection{Identical formal order}

\begin{theorem}[Equivalence of polynomial reproduction; \New]
\label{thm:mls}
Let $\bm A_i$ be non-singular. Then for every polynomial $P$ of total degree at most
$m$,
\begin{equation*}
\sum_{j\in\mathcal N_i}\Phi^{q}_{ji}\,P(\bm r_j)=P(\bm r_i+\bm\delta_q),
\end{equation*}
exactly. Consequently the operator \eqref{eq:mls-weights} has the same leading
truncation error $O(h_i^{\,m+1})$ as the operator of Definition~\ref{def:interp}.
\end{theorem}

\begin{proof}
Since $\bm q$ spans $\mathbb P_m$ in the scaled variable, there exists $\bm c$ with
$P(\bm r_j)=\bm q_{ji}^{\mathsf T}\bm c$ for all $j$, and
$P(\bm r_i+\bm\delta_q)=\bm q(\bm S_i\bm\delta_q)^{\mathsf T}\bm c$. Then
\begin{equation*}
\sum_j\Phi^{q}_{ji}P(\bm r_j)
=\bm q(\bm S_i\bm\delta_q)^{\mathsf T}\bm A_i^{-1}
\sum_j\theta_{ji}\bm q_{ji}\bm q_{ji}^{\mathsf T}\bm c
=\bm q(\bm S_i\bm\delta_q)^{\mathsf T}\bm A_i^{-1}\bm A_i\bm c
=\bm q(\bm S_i\bm\delta_q)^{\mathsf T}\bm c . \qedhere
\end{equation*}
\end{proof}

\begin{center}
\fbox{\parbox{0.95\textwidth}{
\textbf{Consequence for the contribution statement.} Theorem~\ref{thm:mls} shows that
polynomial consistency of degree $m$, and hence interpolation of order $m+1$, is a
property of the consistency constraint and not of the LABFM basis. Any claim that
high-order consistency on a node cloud is specific to LABFM is therefore unfounded and
must not be made. What existing meshless lattice Boltzmann formulations lack is not
the capability but the realisation: the highest polynomial augmentation reported
in~\cite{Strzelczyk2023} is degree four, in two dimensions.
}}
\end{center}

\subsection{The single structural difference: conditioning}

Stack the rows $\bm q_{ji}^{\mathsf T}$ into $\bm Q\in\mathbb R^{N_i\times P}$ and let
$\bm\Theta=\operatorname{diag}(\theta_{ji})$. Then
\begin{equation}
\bm A_i=\bm Q^{\mathsf T}\bm\Theta\bm Q
=\bigl(\bm\Theta^{1/2}\bm Q\bigr)^{\mathsf T}\bigl(\bm\Theta^{1/2}\bm Q\bigr),
\qquad\text{whence}\qquad
\kappa_2(\bm A_i)=\kappa_2\!\bigl(\bm\Theta^{1/2}\bm Q\bigr)^{2}.
\label{eq:gram}
\end{equation}
The moving least squares system is a Gram matrix, formed by normal equations, and its
condition number is therefore the \emph{square} of that of the underlying weighted
design matrix. By contrast, stacking $\widetilde{\bm X}_{ji}^{\mathsf T}$ into
$\bm X\in\mathbb R^{N_i\times p}$ and $\bm W_{ji}^{\mathsf T}$ into $\bm W$,
\begin{equation}
\bm M_i=\bm X^{\mathsf T}\bm W ,
\label{eq:petrov}
\end{equation}
a product of two \emph{distinct} matrices --- a Petrov--Galerkin rather than a Galerkin
construction --- which is not a Gram matrix and involves no squaring. In the language of
\S\ref{sec:structure} the two arms are the two selectors of \eqref{eq:selector}, acting
on one and the same constraint set \eqref{eq:underdet}; by
Proposition~\ref{prop:degenerate} they even coincide when $N_i=p$, so the entire
difference between them lies in how the surplus $N_i-p$ degrees of freedom are spent. The Hermite
ABFs~\eqref{eq:hermiteABF} are moreover orthogonal with respect to a Gaussian weight
while $\bm X$ consists of monomials, so that $\bm M_i$ is close to block diagonal for
near-uniform clouds; the corresponding structure, including entries that vanish to
machine precision by parity, is displayed in~\cite[Fig.~5]{KLN2020}.

\begin{center}
\fbox{\parbox{0.95\textwidth}{
\New\ \textbf{Hypothesis H1 (to be tested at G1).} On the production node cloud and at
matched $m$, $h/s$, kernel and scaling,
$\kappa(\bm A_i)\sim\kappa(\bm M_i)^{2}$, with the gap widening as $m$ increases; and
the largest admissible $m$ is correspondingly larger for the LABFM arm.

\smallskip
\textbf{H1 is the only structural discriminator between the two arms that survives
Theorem~\ref{thm:mls}.} If it is not supported by measurement, there is no theoretical
ground for preferring LABFM over moving least squares in this application.
}}
\end{center}

\subsection{Properties that are not discriminators}

The following are shared by both arms and must not be advanced as reasons to prefer
either:
\begin{enumerate}[label=(\alph*)]
\item formal order of accuracy (Theorem~\ref{thm:mls});
\item a single factorisation per node serving all nineteen departure points
(Remark~\ref{rem:shared} for LABFM; the node-centred kernel for MLS);
\item automatic one-sided operators where the support is incomplete, provided the
local matrix remains non-singular;
\item exact preservation of constants (Corollary~\ref{cor:const} by the difference
form for LABFM; the partition of unity for MLS);
\item applicability of the diagnostics $\Lambda^{q}_i$ and $\max_{\bm k}|G_q(\bm k)|$.
\end{enumerate}
Preprocessing cost in fact favours the control arm: $\bm A_i$ is symmetric positive
definite and admits a Cholesky factorisation at $O(P^{3}/3)$, whereas $\bm M_i$ is
non-symmetric and requires $O(2p^{3}/3)$.

\subsection{Fairness conditions}
\label{sec:fairness}

The comparison is invalid unless all of the following hold. They are stated here so
that a failure to satisfy them is detectable in review.
\begin{enumerate}[label=(F\arabic*)]
\item Identical node cloud, identical stencil membership, identical support radius
$2h_i$ and identical $h/s$.
\item Both arms use the local affine rescaling $\bm S_i$ of \S\ref{sec:affine}.
\item Both arms use a shifted and scaled polynomial basis: the row preconditioning
\eqref{eq:precond} for LABFM, and the scaled argument $\bm q(\bm S_i\rji)$ for MLS.
Comparing a preconditioned LABFM system against an unscaled MLS system is a straw man.
\item Identical kernel, $\psi=\theta=$ Wendland C2.
\item Identical diagnostics, computed by the same code path: $\kappa$, the residual of
the local solve, $\Lambda^{q}_i$, $\max_{\bm k}|G_q(\bm k)|$ at $n_k=16$, and the
monomial reproduction residual of Theorem~\ref{thm:exact}
respectively Theorem~\ref{thm:mls}.
\item Identical arithmetic precision and identical factorisation library.
\end{enumerate}

\subsection{Decision rule}
\label{sec:decision}

At gate G1 both arms are run on the production node cloud and the following are
reported for each: the largest $m$ admitting $\kappa<10^{10}$; $\kappa(m)$ at the
target order; $\Lambda^{q}_i$ and $\max_{\bm k}|G_q|$ distributions; the factorisation
cost; and the stencil radius in index units. Then
\begin{itemize}
\item if the LABFM arm admits a larger $m$, or exhibits a condition-number advantage of
at least two orders at the target $m$, adopt LABFM; H1 is supported;
\item if the two arms are comparable, adopt the control arm --- it is simpler, cheaper
to factorise and more widely implemented --- and drop the LABFM designation;
\item if the control arm is superior, adopt it.
\end{itemize}
In all three outcomes the theory of \S\ref{sec:interp} is unaffected: Theorems
\ref{thm:exact} and \ref{thm:affine}, Corollaries \ref{cor:const} and \ref{cor:zero},
the amplification constant \eqref{eq:lebesgue} and the stability criterion
\eqref{eq:vn-crit} hold verbatim for either construction, and the contribution
statement of \S1.3 is method-agnostic.

%========================================================================
\section{Verification ladder}
\label{sec:ladder}

Each rung reproduces either a published result or an identity proved above. A rung must
pass before the next is attempted.

\begin{widetable}
\begin{longtable}{@{}P{1.05cm}P{3.3cm}P{11.35cm}@{}}
\toprule
\textbf{Rung} & \textbf{Target} & \textbf{Pass condition} \\
\midrule
\endfirsthead
\toprule
\textbf{Rung} & \textbf{Target} & \textbf{Pass condition} \quad\emph{(cont.)} \\
\midrule
\endhead
\midrule
\multicolumn{3}{r@{}}{\footnotesize\emph{continued on next page}}\\
\endfoot
\bottomrule
\endlastfoot
L0 & \cite[Eqs.~(15),(16)]{KLN2020} & Basis-dimension generator returns
$\{5,14,27,44\}$ in two dimensions for $m=2,4,6,8$ and $\{3,9,19,34,55\}$ in three
dimensions for $m=1,\dots,5$. \\
L0b & Eq.~\eqref{eq:parity} & On a uniform Cartesian cloud the parity defect
\eqref{eq:parity-defect} satisfies $\eta_i\lesssim10^{-14}$. Cheapest available test of
the multi-index ordering and the Hermite recurrence; requires no solve. \\
L1 & \cite[\S 4.1]{KLN2020} & Two dimensions, original conic ABFs, distribution noise
$\varepsilon/s=0.5$, $h/s=2$, test function \cite[Eq.~(23)]{KLN2020}: gradient
converges at order $m$ and the Laplacian at order $m-1$ for $m=2,4,6$.
\textbf{The truncation of Remark~\ref{rem:truncation} must be applied} or the numbers
will not match. \\
L2 & \cite[\S 3.2]{KLN2020} & Two-dimensional critical ratios recovered with the
\emph{original} ABFs: $(h/s)_{\rm crit}\approx\{0.75,1.15,1.60,2.10\}$ and
$N_{\rm crit}\approx\{8,21,37,57\}$. Establishes the measurement procedure later
applied to the Hermite construction (open items U1, U8). \\
L3 & \cite[\S 3]{KL2021} & \emph{Negative test.} Three dimensions, original ABFs,
$m=4$: $\bm M$ of size $34$ with rank $33$, independently of the node distribution. A
full-rank result indicates an error in the ABF assembly, not an improvement. \\
L4 & \cite[\S 3, Fig.~2]{KL2021} & Hermite ABFs reduce $\kappa(\bm M)$ by several
orders of magnitude relative to the original ABFs, and restore full rank in three
dimensions at $m=4$. \\
L5a & Corollary~\ref{cor:zero} & $\bm\delta_q=\bm0$ yields $w^{q}\equiv0$ and an
unchanged field to machine precision. \\
L5b & Theorem~\ref{thm:exact} & All monomials of total degree at most $m$ reproduced at
the departure point to a residual below $10^{-12}$. \\
L5c & Corollary~\ref{cor:const} & A uniform field is unchanged after streaming in all
nineteen directions. \\
L5d & Theorem~\ref{thm:mls} & \emph{Control arm.} The moving least squares operator
\eqref{eq:mls-weights} reproduces all monomials of total degree at most $m$ at the
departure point to a residual below $10^{-12}$, and satisfies
$\sum_j\Phi^{q}_{ji}=1$ to machine precision. \\
L5e & Proposition~\ref{prop:degenerate} & With $N_i=p$ the weights equal
$\bm X^{-\mathsf T}\bm C$ and are independent of the selector: both arms, and any
choice of ABF, must agree bitwise, and in one dimension must reproduce the Lagrange
weights on $m+1$ points. Validates the assembly without exercising the basis at all. \\
L5f & Proposition~\ref{prop:gauge} & \emph{Gauge test.} Rebuilding $\bm X$ and $\bm C$
about a displaced origin $\bm r_0=\bm r_i+\varepsilon\bm n$, $\varepsilon\sim h$,
with $\bm S$ and $\mathcal N_i$ unchanged, must leave $\bm w$ invariant to a relative
$10^{-10}$. Isolates errors in the monomial assembly and multi-index ordering, which
rung L5b detects only at the highest degrees. \\
L6 & \cite[\S 4.2, Table 1]{KLN2020} & Trend of the Laplacian stability limit with $h/s$ and
$m$ recovered; corroborates the ABF implementation. \\
L7 & Eq.~\eqref{eq:vn-crit} & On the production node cloud,
$\max_{\bm k}|G_q(\bm k)|\le1.01$ for all nodes and all $q$. \\
L7b & \S\S\ref{sec:project-stencil}, \ref{sec:boundary} &
\emph{Boundary admissibility.} On every one-sided and nearly one-sided P6/31 cloud,
report the three numerical ranks, equilibrated rcond, $\lambda_0$, $T_7$, all final-solve
gates and the direction-specific $\Lambda_q$; then measure
$\max_{\bm k}|G_q(\bm k)|$ for the assembled interpolation operator. Local gates settle
constructibility, while the amplification test settles whether the derivative-instability
result transfers; together they close open item U4. \\
L8 & Hypothesis H1, \S\ref{sec:decision} & \emph{Head-to-head.} Both arms run on the
production node cloud under the fairness conditions of \S\ref{sec:fairness}; the
decision rule of \S\ref{sec:decision} is applied and its outcome recorded. \\
\end{longtable}
\end{widetable}

%========================================================================
\section{Open items}

No item in this table may be used as an implementation assumption until it has been
measured and the result written back into this document.

\begin{widetable}
\begin{longtable}{@{}P{0.8cm}P{3.5cm}P{11.4cm}@{}}
\toprule
\textbf{\#} & \textbf{Item} & \textbf{Why open, and how it is to be closed} \\
\midrule
\endfirsthead
\toprule
\textbf{\#} & \textbf{Item} & \textbf{Why open} \quad\emph{(cont.)} \\
\midrule
\endhead
\midrule
\multicolumn{3}{r@{}}{\footnotesize\emph{continued on next page}}\\
\endfoot
\bottomrule
\endlastfoot
U1 & $(h/s)^{\rm 3D}_{\rm crit}$ & Not in the literature; \cite{KL2021} defers three
dimensions. Measured on the production grid for $m=6$ with a metric-ball support:
$\operatorname{rank}\bm X=75$--$80$ of $83$ at $h/s=1.5$ and $83$ at $h/s\ge1.6$
($N_i=142$), so $(h/s)^{\rm 3D}_{\rm crit}\simeq1.6$ for a lattice cloud. Whether this
constant enters a future full three-dimensional path depends on a selector that is not
implemented. The current two-dimensional P6/31 selector is fully specified in
\S\ref{sec:project-stencil} and does not assume this three-dimensional value. \New \\
U2 & Unisolvency of structured lattice balls & \textbf{Closed, negatively.} The $123$-point radius-three lattice ball at $h/s=1.5$ attains
$\operatorname{rank}\bm X=77$ of $83$; the singular spectrum has a hard gap
($\sigma_{79}/\sigma_1=4.9\times10^{-8}$, $\sigma_{80}/\sigma_1=1.0\times10^{-17}$), so
the deficiency is exact and not a threshold artefact. It occurs at mid-channel and at
the flat upper wall as well as on the hill face, hence it is a property of the lattice
ball and not of the metric. A sufficient design screen is the slicing count: the support
must decompose into $m+1$ parallel slices holding at least
$\dim\mathbb P_{m-\ell}^{2\rm D}=28,21,15,10,6,3,1$ points; this predicts the measured
outcome at $h/s=1.5$, $1.6$ and $1.8$ correctly. \New \\
U3 & Admissible bound on $\Lambda^{q}_i$ & No counterpart exists in the literature,
which does not define an interpolation operator. Reference values are now available
on a $25$-point in-plane support: $1.2$--$2.4$ in the interior, $3.0$ in the wall band; a
rejected configuration reaches $3.5\times10^{4}$ and a ghost-fed one
$1.2\times10^{3}$. The current P6/31 implementation uses a distinct,
direction-independent hard gate $\lambda_0\le512$ in the compact tier and
$\lambda_0\le4096$ in repair, while the actual direction-specific
$\Lambda_i^q>6$ remains a warning rather than rejection. The admissible production
bound on $\Lambda_i^q$ is therefore still open and must be fixed by the error and
stability budget. Close jointly with rung L7 and U5. \\
U4 & \textbf{Near-wall treatment} & \textbf{Partly closed at the local-algebra level.}
The current selector uses no wall ghost: it clamps a physical seven-layer P6/31 cloud
into the domain and fails closed unless the rank, equilibrated-rcond, $\lambda_0$ and
three-solver gates of \S\ref{sec:project-stencil} pass. This addresses numerical
constructibility for the audited grids and avoids treating the adverse result of
\cite[\S 7.1]{KL2021} as a universal impossibility theorem. Still open are the global
interpolation amplification spectrum, observed wall convergence over the supported grid
family, and production physical validation. Close by rung L7b; until then, local gate
success must not be reported as proof of wall-stable seventh-order production accuracy. \\
U5 & Interpolation error against the lattice Boltzmann error floor & Determines whether
raising the interpolation order changes the solution at all, given the saturation law
of~\cite{Strzelczyk2023}. Close by an offline measurement on a converged field. \\
U6 & Measured roofline position & Determines whether a reduction in operation count
converts to wall-clock time. Close by profiling the existing kernel. \\
U7 & Hypothesis H1: which weight construction & Theorem~\ref{thm:mls} removes order of
accuracy as a discriminator, leaving conditioning \eqref{eq:gram} versus
\eqref{eq:petrov} as the only structural difference. Close by rung L8 under the
fairness conditions of \S\ref{sec:fairness}; the decision rule is \S\ref{sec:decision}.
Until U7 is closed, the method designation in all external material is to be
``consistency-based meshless departure-point interpolation'', not ``LABFM''. \\
U8 & Critical ratio for the Hermite ABFs & The published critical values belong to the
\emph{original} ABFs, which the production path does not use; only recommended
operating values with an unquantified margin are published for the Hermite
construction (\S\ref{sec:sampling}). The present P6/31 implementation bypasses a fixed
published ratio by trying finite $h$ and anisotropy ladders and accepting only the
measured matrix gates; that is an implementation policy, not a determination of a
universal critical ratio. Close together with U1, in two dimensions first so that rung
L2 provides a validated procedure. \\
U9 & Whether the extrusion symmetry survives & The weight table presumes a node cloud
that is a uniform extrusion along the homogeneous direction. Relaxing this in Stage~2
multiplies the table by the number of nodes in that direction. Close by deciding, before
Stage~2, whether variable resolution is wanted in the homogeneous direction at all. \
\end{longtable}
\end{widetable}

%========================================================================
\section{Interface with the existing solver}

\begin{itemize}
\item A compile-time switch selects the present streaming path; it is mutually
exclusive with the existing streaming paths, enforced at compile time.
\item No existing streaming, collision or infrastructure module is modified.
\item The collision operator, the forcing scheme, the statistics, the file
input/output, the halo exchange, the mass correction and its monitoring are reused
unchanged.
\item \textbf{Ghost storage is retained, but it is not part of a LABFM cloud.}
The incumbent wall boundary, ITB/GILBM streaming and statistics-window contracts still
require their existing ghost layers and must not be dismantled. Independently, the
P6/31 selector of \S\ref{sec:project-stencil} admits canonical physical IDs only and
clamps its one-sided stencil into the domain. If a finite-difference hybrid in the sense
of \cite[Table 1]{KL2021} is adopted later, the interface between that strip and the
meshless interior must be specified as a separate method change before any production
run is submitted.
\item The stencil is centred on the arrival node, which restricts the time step through
Proposition~\ref{prop:cfl}. Any change to the Courant number must be checked against
both that bound and the relaxation-time bound.
\item The scheme is specified for a single global time step. Local time stepping is
out of scope.
\end{itemize}

%========================================================================
\appendix
\section{Provenance table}
\label{app:prov}

\begin{widetable}
\begin{longtable}{@{}P{3.9cm}P{4.45cm}P{7.3cm}@{}}
\toprule
\textbf{This report} & \textbf{Source} & \textbf{Content} \\
\midrule
\endfirsthead
\toprule
\textbf{This report} & \textbf{Source} & \textbf{Content} \quad\emph{(cont.)} \\
\midrule
\endhead
\midrule
\multicolumn{3}{r@{}}{\footnotesize\emph{continued on next page}}\\
\endfoot
\bottomrule
\endlastfoot
\eqref{eq:X}, \eqref{eq:D} & \cite[Eqs.~(1),(2)]{KLN2020}, \cite[Eqs.~(1),(2)]{KL2021}
& Monomial and derivative vectors \\
\eqref{eq:taylor} & \cite[Eq.~(3)]{KLN2020} & Compact Taylor expansion \\
\eqref{eq:Lop} & \cite[Eq.~(4)]{KLN2020}, \cite[Eq.~(3)]{KL2021} & General discrete
operator; $\bm C^{d}$ from \cite[Eq.~(5)]{KLN2020}, \cite[Eq.~(4)]{KL2021} \\
\eqref{eq:wPsi} & \cite[Eq.~(10)]{KLN2020}, \cite[Eq.~(5)]{KL2021} & ABF expansion of
the weights \\
\eqref{eq:moments} & \cite[Eqs.~(13),(14)]{KLN2020}, \cite[Eqs.~(6),(7)]{KL2021} &
Moment matrix and local system \\
\eqref{eq:p} & \cite[Eqs.~(15),(16)]{KLN2020} & Basis dimension in two and three
dimensions \\
\eqref{eq:origABF} & \cite[Eq.~(19)]{KLN2020} & Original ABFs \\
\eqref{eq:hermiteABF} & \cite[Eqs.~(10),(11)]{KL2021} & Hermite ABFs \\
\eqref{eq:precond} & \cite[Eq.~(9)]{KL2021} & Row preconditioning by powers of $h$ \\
\S\ref{sec:sampling} table & \cite[\S 3.2]{KLN2020}, \cite[\S\S 4, 5]{KL2021} &
Critical (original ABFs) and recommended (Hermite) stencil ratios; stencil optimisation \\
\S\ref{sec:boundary} & \cite[\S 4.1]{KLN2020}, \cite[\S 7.1, Table 1,
Eq.~(24), Appendix A]{KL2021} & Boundary treatment; the 2021 retraction of the 2020
one-sided claim and the finite-difference hybrid \\
\midrule
\S\ref{sec:project-stencil},
\eqref{eq:mandatory-window}--\eqref{eq:T7-selector} & Current project sources:
\texttt{labfm\_stencil.h}, \texttt{labfm\_solve.h} & Two-dimensional P6/31-point
candidate retention, mandatory template/QRCP membership, anisotropic scaling, rank and
rcond gates, and compact-first repair selection; an implementation record, not a claim
from the cited LABFM literature \\
\eqref{eq:N3D} & \New & Three-dimensional neighbour count; geometric extension of the
two-dimensional relation in \cite[\S 2]{KL2021} \\
\eqref{eq:departure}--\eqref{eq:stream} & \New & Departure-point interpolation
operator (Definition~\ref{def:interp}) \\
Theorem~\ref{thm:exact} & \New & Exact polynomial reproduction at the departure point \\
Corollaries~\ref{cor:const}, \ref{cor:zero} & \New & Constant preservation; zero
displacement identity \\
Remark~\ref{rem:noninterp} & \New & Approximating, not interpolating; invalidates a
natural but incorrect regression test \\
\eqref{eq:lebesgue} & \New & Amplification constant of the interpolation operator \\
\eqref{eq:amp}, \eqref{eq:vn-crit} & \New, adapted from
\cite[\S 4, Eqs.~(13a)--(13c)]{KL2021} & Von Neumann amplification and stability
criterion for the streaming step \\
Theorem~\ref{thm:affine}, \eqref{eq:hermite3d} & \New & Affine invariance; three
dimensional Hermite ABFs \\
Remark~\ref{rem:distribution}, \eqref{eq:measure} & \New & Distributional reading;
$\bm C$ as the moment vector of a target distribution \\
Remark~\ref{rem:truncation} & \cite[Appendix A]{KLN2020} & Truncation of the original
ABFs for $m\ge4$, required to reproduce the published results \\
Remark~\ref{rem:innerprod}, \eqref{eq:sampling-norm} & \New & Normalisation of the
sampling condition \\
\eqref{eq:parity}, \eqref{eq:parity-defect} & \New & Parity structure of the continuum
limit; the parity defect as a leading sampling indicator \\
\eqref{eq:underdet}, \eqref{eq:selector} & \New & Underdetermination of the weight
system; the selector formulation unifying both arms \\
Propositions~\ref{prop:degenerate}, \ref{prop:gauge}, Remarks~\ref{rem:surplus},
\ref{rem:origins} & \New & Degenerate limit $N_i=p$; gauge freedom of the expansion
origin and the two origins that are not free \\
Remark~\ref{rem:threepoints}, Proposition~\ref{prop:cfl}, \eqref{eq:cfl-bound} & \New &
Roles of $\bm r_i$, $\bm r_j$, $\bm r^{*}_q$; Courant bound for a node-centred stencil \\
\eqref{eq:mls} & \cite{Lancaster1981,Mirzaei2012,Trask2016} & Moving least squares;
node-centred (generalised) kernel variant \\
\eqref{eq:mls-weights}, Theorem~\ref{thm:mls} & \New & Departure-point form of the
control operator; equivalence of polynomial reproduction \\
\eqref{eq:gram}, \eqref{eq:petrov}, H1 & \New & Gram versus Petrov--Galerkin
conditioning; the sole surviving discriminator \\
\S\S\ref{sec:fairness}, \ref{sec:decision} & \New & Fairness conditions and decision
rule for the head-to-head comparison \\
\end{longtable}
\end{widetable}

%========================================================================
\section{A worked reading of \texorpdfstring{\cite[\S 3.1]{KLN2020}}{King et al. Sec. 3.1}}
\label{app:tutorial}

The construction of \cite[\S 3.1]{KLN2020} is short, and its brevity conceals what is
being asserted. This appendix restates it in full, and then carries it out in closed
form on the smallest non-trivial configuration, so that every claim can be checked by
hand.

\begin{center}
\fbox{\parbox{0.95\textwidth}{
\textbf{Precedence.} This appendix is a restatement, not a source. Sections
\ref{app:statement}--\ref{app:dims} assert nothing that is not in
\cite[\S\S 2, 3.1]{KLN2020}; their correctness consists in being a faithful
transcription, so they cannot serve as the authority for their own content.
\textbf{Where this appendix and \cite{KLN2020} disagree, \cite{KLN2020} governs and
this appendix is in error.} Only the worked example of \S\S\ref{app:example},
\ref{app:under}, Proposition~\ref{prop:dims} and the material marked \New{} stand on
their own, and they function as checks on the transcription rather than substitutes for
the source. The reader is expected to return to \cite{KLN2020} whenever rungs L1--L4 or
L6 fail, whenever a parameter such as $h/s$ or $m$ is to be chosen, whenever the work is
to be cited externally, and --- above all --- whenever a statement here is doubted.
}}
\end{center} The example also exhibits, in a setting where all quantities are explicit, three
properties established abstractly in the main text: the cancellation of the basis
functions in the degenerate case (Proposition~\ref{prop:degenerate}), the parity
structure of the moment matrix \eqref{eq:parity}, and the necessity of a lower bound on
$h/s$ (\S\ref{sec:sampling}).

\subsection{The statement to be proved}
\label{app:statement}

Section~2 of \cite{KLN2020} establishes an identity rather than an approximation.
Substituting the truncated Taylor expansion \eqref{eq:taylor} into the discrete operator
\eqref{eq:Lop} gives \cite[Eq.~(7)]{KLN2020}
\begin{equation}
\mathcal L^{d}_i(\phi)
=\underbrace{\Bigl(\sum_j x_{ji}w^{d}_{ji}\Bigr)}_{B^{d,1}_x}\frac{\partial\phi}{\partial x}
+\underbrace{\Bigl(\sum_j y_{ji}w^{d}_{ji}\Bigr)}_{B^{d,1}_y}\frac{\partial\phi}{\partial y}
+\underbrace{\Bigl(\sum_j \tfrac{x_{ji}^{2}}{2}w^{d}_{ji}\Bigr)}_{B^{d,2}_{xx}}
\frac{\partial^{2}\phi}{\partial x^{2}}+\cdots
\label{eq:momexpand}
\end{equation}
The discrete operator is therefore \emph{exactly} a linear combination of the
derivatives at node $i$, in which the coefficient of each derivative is a moment of the
weights. Approximating a prescribed derivative is consequently not a question of
choosing a good formula but of solving an algebraic problem: set the moment belonging
to the desired derivative to unity and every other moment to zero. The moments that are
required to vanish are precisely the error terms of the title of
\cite[\S 3.1]{KLN2020}; this is the sense in which they are ``eliminated'', and the
elimination is exact, not asymptotic.

Smoothed particle hydrodynamics obtains its weights from a kernel derivative
\cite[Eq.~(9)]{KLN2020} and therefore accepts whatever moments that kernel happens to
possess; \cite{KLN2020} observes that the result ``is not particularly useful'', being
of zero order. The contribution of \cite[\S 3.1]{KLN2020} is a constructive means of
imposing the moments instead of inheriting them.

\subsection{The construction}

The obstruction is that the moment conditions are $p$ equations in $N_i$ unknowns and,
as discussed in \S\ref{sec:structure}, are underdetermined whenever $N_i>p$. The device
of \cite[\S 3.1]{KLN2020} is to \emph{parameterise} the weights rather than solve for
them: by \eqref{eq:wPsi} the $N_i$ unknown weights are replaced by $p$ unknown
coefficients multiplying $p$ prescribed functions of the separation. Substituting this
parameterisation into the moment vector and extracting the coefficients from the sum
produces a square system, \eqref{eq:moments}. The two forms of that system carry
opposite intent and should not be conflated: \cite[Eq.~(12)]{KLN2020} maps a given
$\bm\Psi$ to the moments it produces, whereas \cite[Eq.~(14)]{KLN2020} prescribes the
moments and solves for $\bm\Psi$. Only the second is the method.

\subsection{What the dimensions are free to be, and what fixes them}
\label{app:dims}

\cite[\S 3.1]{KLN2020} closes the construction with a sentence that is easy to read as
bookkeeping:

\begin{quote}
``The above analysis imposes no constraint on the lengths of $X$, $W$, $C$, and the size
of $M$, although for consistency between (13) and (14), we require $X$, $W$ and $C$ to
be of equal length, and $M$ square.''
\end{quote}

\noindent It is not bookkeeping. Read carefully it separates a genuine design decision
from two consequences of it, and it is worth unpacking because the freedom it does
\emph{not} remove is the freedom the rest of this report depends on.

\paragraph{The first clause is literally true.} Nothing in
\eqref{eq:taylor}--\eqref{eq:moments} refers to a dimension. The moment expansion
\eqref{eq:momexpand} has infinitely many terms; $\bm X$ and $\bm W$ are formally
infinite sequences; the outer product is formally an infinite array. Truncation is
\emph{imposed} at this point in the argument, not derived from anything preceding it.
That is the content of ``imposes no constraint''.

\paragraph{Of the three stated equalities, only one is a choice.} Write
$n_X=|\bm X|$, $n_W=|\bm W|$, $n_C=|\bm C|$. Then $\bm M$ has shape $n_X\times n_W$,
so $\bm M\bm\Psi$ lies in $\mathbb R^{n_X}$, and \eqref{eq:moments} equates it to
$\bm C$. Hence
\begin{equation}
n_C=n_X\quad\text{is forced by the equation being well formed,}
\label{eq:nCnX}
\end{equation}
and is not a modelling decision at all. The only genuine choice is $n_W$ against $n_X$,
and setting $n_W=n_X\equiv p$ is what makes $\bm M$ square. The three lengths carry
three different meanings, and it is their coincidence that is the design decision:
\begin{center}
\begin{tabular}{@{}lll@{}}
\toprule
Vector & Its length counts & Status \\
\midrule
$\bm X$ & the moment conditions imposed, i.e. the error terms annihilated & sets $p$ \\
$\bm W$ & the free coefficients available in $\bm\Psi$ & \textbf{the choice} \\
$\bm C$ & the values prescribed & forced, \eqref{eq:nCnX} \\
\bottomrule
\end{tabular}
\end{center}
\noindent In words: \emph{one free coefficient per moment condition}. The squareness of
$\bm M$ is the consequence, not the motive.

\begin{proposition}[The two alternatives; \New]
\label{prop:dims}
Let $n_X=p$ and let $\bm M$ have full rank. Then
\begin{enumerate}[label=(\roman*)]
\item if $n_W>p$, the system \eqref{eq:moments} is underdetermined, its solution set is
an affine subspace of dimension $n_W-p$, and \emph{every} member satisfies all $p$
moment conditions; polynomial reproduction, and hence the order, is unaffected, but a
further rule is required to select $\bm\Psi$;
\item if $n_W<p$, the system is overdetermined and generically inconsistent; no
$\bm\Psi$ satisfies all $p$ conditions, the unsatisfied moments survive in
\eqref{eq:momexpand}, and the order falls.
\end{enumerate}
\end{proposition}

\begin{proof}
Immediate from the shape $p\times n_W$ of $\bm M$ and the rank hypothesis. For (i),
membership of the solution set means $\bm M\bm\Psi=\bm C$ holds exactly, which is the
definition of the moment conditions; Theorem~\ref{thm:exact} then applies verbatim. For
(ii), a least-squares surrogate leaves a non-zero residual in at least one moment, and
that moment multiplies its derivative in \eqref{eq:momexpand}.
\end{proof}

\begin{remark}[$n_W$ interpolates between two extremes]
\label{rem:nW}
Case (i) is not pathological; it is a family of admissible schemes that
\cite{KLN2020} declines to enter. Its extreme is instructive. Since
$\bm w=\bm W\bm\Psi$, the reachable weights lie in the column space of $\bm W$, of
dimension $n_W$. At $n_W=p$ the ansatz determines $\bm w$ completely once $\bm C$ is
given; at the opposite extreme $n_W=N_i$ with $\bm W$ invertible, the reachable set is
all of $\mathbb R^{N_i}$ and the ansatz imposes \emph{nothing}, returning the bare
underdetermined system \eqref{eq:underdet}. The choice $n_W=p$ therefore places LABFM at
the minimal end of that range: the smallest parameterisation that still closes the
system. Every intermediate value would introduce a second layer of arbitrariness on top
of the selection already discussed in \S\ref{sec:structure}.
\end{remark}

\paragraph{Square is necessary, not sufficient.} The quoted sentence secures only the
shape of $\bm M$. Its invertibility is a separate matter and is where the method
actually fails in practice: through the sampling condition of \S\ref{sec:sampling}, and,
in three dimensions with the original basis, through the unconditional rank deficiency
of \S\ref{sec:abf}.

\paragraph{A second restriction, which the sentence does not state.} Equality of the
three lengths does not by itself determine $p$; \cite[Eqs.~(15),(16)]{KLN2020} do, and
they admit only the values $2,5,9,14,\dots$ in two dimensions and $3,9,19,34,55,83$ in
three. The reason is that $p$ must fall at the end of a complete degree. Truncating in
mid-degree --- taking, in two dimensions, $\bm X=[x,\,y,\,x^{2}/2]$ with $p=3$ ---
controls $B_{x}$, $B_{y}$ and $B_{xx}$ while leaving $B_{xy}$ and $B_{yy}$ free. Since
$B^{(n)}\sim h^{\,n-l}$, those surviving second moments contribute at $O(h)$ to a first
derivative where a complete truncation would have left $O(h^{2})$: the order drops, and
it drops \emph{anisotropically}, a field varying along $x$ retaining an accuracy that a
field varying along $y$ does not. Equivalently, $\operatorname{span}\{1,x,y,x^{2}\}$ is
not invariant under rotation of the axes whereas $\mathbb P_2$ is, so the operator would
depend on an orientation that is arbitrary on a body-fitted curvilinear grid. This is
enforced structurally in the implementation of \S\ref{sec:abf}: the multi-index
enumeration loops over complete degrees, so a mid-degree truncation cannot be expressed.

\paragraph{What remains free.} The sentence locks the ansatz dimension to the constraint
count. It says nothing about $N_i$, and $N_i$ is not locked: the necessary condition is
only $N_i>p$, and the surplus $N_i-p$ is precisely the freedom exploited in
\S\ref{sec:structure} and contested between the two arms of \S\ref{sec:mls}. The
division is therefore sharp, and worth stating as the summary of this subsection:
\emph{the number of unknowns in the parameterisation is fixed by the number of
conditions; the number of data points is not.}

\subsection{A minimal example}
\label{app:example}

Let the problem be one-dimensional with $m=2$, so that
$\bm X_{ji}=\bigl[x_{ji},\,x_{ji}^{2}/2\bigr]^{\mathsf T}$ and $p=m=2$; let node $i$ lie
at $x=0$ and let its neighbours lie at $x=\pm s$. Then
\begin{equation}
\bm X_{j_1 i}=\Bigl[s,\ \tfrac{s^{2}}{2}\Bigr]^{\mathsf T},
\qquad
\bm X_{j_2 i}=\Bigl[-s,\ \tfrac{s^{2}}{2}\Bigr]^{\mathsf T},
\qquad
\det\bm X=s^{3}\neq0 ,
\label{eq:exX}
\end{equation}
so the moment conditions $\bm X^{\mathsf T}\bm w=\bm C$ possess a unique solution for
every right-hand side. Written out, they are
\begin{equation}
s\,(w_1-w_2)=C_1,
\qquad
\tfrac{s^{2}}{2}\,(w_1+w_2)=C_2 .
\label{eq:excond}
\end{equation}

\paragraph{First derivative.} Setting $\bm C=[1,0]^{\mathsf T}$, the second of
\eqref{eq:excond} forces $w_2=-w_1$ --- this is the elimination of the $\phi''$ error
term --- whence $w_1=1/2s$ and $w_2=-1/2s$, and
\begin{equation}
\mathcal L^{x}_i(\phi)=\frac{\phi_{j_1}-\phi_i}{2s}-\frac{\phi_{j_2}-\phi_i}{2s}
=\frac{\phi(s)-\phi(-s)}{2s}.
\label{eq:excd1}
\end{equation}

\paragraph{Second derivative.} Setting $\bm C=[0,1]^{\mathsf T}$, the first of
\eqref{eq:excond} forces $w_1=w_2$ --- the elimination of the $\phi'$ error term ---
whence $w_1=w_2=1/s^{2}$, and
\begin{equation}
\mathcal L^{L}_i(\phi)=\frac{\phi_{j_1}-\phi_i}{s^{2}}+\frac{\phi_{j_2}-\phi_i}{s^{2}}
=\frac{\phi(s)-2\phi(0)+\phi(-s)}{s^{2}}.
\label{eq:excd2}
\end{equation}

The machinery of \cite[\S 3.1]{KLN2020} thus returns the classical central and second
differences. This is the expected outcome and is worth stating explicitly, because it
identifies what the method generalises: LABFM is a procedure for constructing, on an
arbitrary node distribution, the object that the moment conditions define uniquely on a
uniform one.

\subsection{The superscript \texorpdfstring{$d$}{d}: what it does and does not affect}
\label{app:super-d}

The label $d$ appears on $w^{d}_{ji}$, $\bm\Psi^{d}_i$, $\bm B^{d}_i$ and $\bm C^{d}$
but not on $\bm X_{ji}$, $\bm W_{ji}$ or $\bm M_i$. That asymmetry is the whole content
of the notation, and it is worth making explicit, because the extension of
\S\ref{sec:interp} consists of nothing but replacing $d$ by a different label.

\paragraph{Dependency chain.} Reading \eqref{eq:moments} and \eqref{eq:wPsi} from right
to left,
\begin{equation}
\bm C^{d}\ \longmapsto\ \bm\Psi^{d}_i=\bm M_i^{-1}\bm C^{d}
\ \longmapsto\ w^{d}_{ji}=\bm W_{ji}\cdot\bm\Psi^{d}_i
\ \longmapsto\ \bm B^{d}_i=\sum_j\bm X_{ji}w^{d}_{ji}\;\bigl(=\bm C^{d}\bigr).
\label{eq:dchain}
\end{equation}
Only $\bm C^{d}$ is supplied; the remaining three are consequences. The superscript is
therefore \emph{passive}: it enters no formula except as the argument of the solve, and
it is carried on $\bm\Psi$, $w$ and $\bm B$ only because those objects inherit the
dependence. In particular $\bm B^{d}_i$ is the vector of moments \emph{achieved}, and
after the solve it coincides with the vector \emph{prescribed}; \cite{KLN2020} writes
$\bm B$ in Eq.~(8), before the ansatz, when it is whatever the weights happen to give,
and replaces it by $\bm C$ in Eq.~(14), when it is demanded.

The passivity is not a redundancy, however: for each $d$ the objects
$\bm\Psi^{d}_i$ and $w^{d}_{ji}$ are genuinely different arrays of numbers, as the
example below shows.

\paragraph{A worked example in two dimensions.} Take $m=2$, so that $p=5$ and
\begin{equation*}
\bm X_{ji}=\Bigl[x_{ji},\ y_{ji},\ \tfrac{x_{ji}^{2}}{2},\ x_{ji}y_{ji},\
\tfrac{y_{ji}^{2}}{2}\Bigr]^{\mathsf T},
\end{equation*}
and place five neighbours at
\begin{equation}
(\pm s,0),\quad (0,\pm s),\quad (s,s),
\label{eq:2dstencil}
\end{equation}
numbered $j_1,\dots,j_5$ in that order, so that $N_i=p$ and, by
Proposition~\ref{prop:degenerate}, the basis functions cancel and the weights are those
of the moment conditions alone. Writing out $\bm X^{\mathsf T}\bm w=\bm C$ row by row,
\begin{align}
s\,(w_1-w_2)+s\,w_5&=C_1, &
s\,(w_3-w_4)+s\,w_5&=C_2, &
\tfrac{s^{2}}{2}(w_1+w_2+w_5)&=C_3,\notag\\
s^{2}w_5&=C_4, &
\tfrac{s^{2}}{2}(w_3+w_4+w_5)&=C_5. & &
\label{eq:2dsys}
\end{align}
Solving for the three standard right-hand sides of \cite[Eq.~(5)]{KLN2020}:

\begin{center}
\begin{tabular}{@{}llll@{}}
\toprule
$d$ & $\bm C^{d}$ & $\bm w^{d}=(w_1,\dots,w_5)$ & resulting operator \\
\midrule
$x$ & $[1,0,0,0,0]$ & $\bigl[\tfrac{1}{2s},\,-\tfrac{1}{2s},\,0,\,0,\,0\bigr]$
& $\dfrac{\phi_E-\phi_W}{2s}$ \\[0.5em]
$y$ & $[0,1,0,0,0]$ & $\bigl[0,\,0,\,\tfrac{1}{2s},\,-\tfrac{1}{2s},\,0\bigr]$
& $\dfrac{\phi_N-\phi_S}{2s}$ \\[0.5em]
$L$ & $[0,0,1,0,1]$ & $\bigl[\tfrac{1}{s^{2}},\,\tfrac{1}{s^{2}},\,\tfrac{1}{s^{2}},
\,\tfrac{1}{s^{2}},\,0\bigr]$
& $\dfrac{\phi_E+\phi_W+\phi_N+\phi_S-4\phi_i}{s^{2}}$ \\
\bottomrule
\end{tabular}
\end{center}

\noindent The three classical formulae are recovered, from one and the same $\bm M_i$,
by changing $\bm C^{d}$ alone. Three features of the table answer the question of what
the superscript does.

\begin{enumerate}[label=(\roman*)]
\item \emph{The weights are different arrays, not a relabelling.} $\bm w^{x}$ is
supported on the two horizontal neighbours, $\bm w^{y}$ on the two vertical ones, and
$\bm w^{L}$ on all four; the fifth neighbour at $(s,s)$ receives zero weight in all
three cases, since $C_4=0$ throughout and \eqref{eq:2dsys} gives $w_5=C_4/s^{2}$.
\item \emph{The superscript carries a dimension.} $w^{x}\sim s^{-1}$ while
$w^{L}\sim s^{-2}$. This is the scaling requirement of \cite[\S 2]{KLN2020}, that
$w^{d}$ scale as $h^{-l}$ with $l$ the order of the derivative; for the interpolation
operator of \S\ref{sec:interp}, $l=0$ and the weights are dimensionless, which is why
the affine rescaling of \S\ref{sec:affine} requires no inverse transformation.
\item \emph{The dependence on $\bm C^{d}$ is linear.} Since
$\bm\Psi^{d}_i=\bm M_i^{-1}\bm C^{d}$ and $w^{d}=\bm W\bm\Psi^{d}$, the map
$\bm C^{d}\mapsto\bm w^{d}$ is linear. Hence operators superpose: with
$\bm C^{xx}=[0,0,1,0,0]$ and $\bm C^{yy}=[0,0,0,0,1]$ one has
$\bm C^{L}=\bm C^{xx}+\bm C^{yy}$ and therefore
$\bm w^{L}=\bm w^{xx}+\bm w^{yy}$, which the table verifies term by term. The
hyperviscosity operators of \cite[Eq.~(25)]{KLN2020} are constructed in exactly this
way, and so is any other linear combination of derivatives one may require.
\end{enumerate}

\begin{remark}[Why this is the entry point for the present work]
\label{rem:d-to-q}
Since $d$ enters only through $\bm C^{d}$, and $\bm M_i$ is indifferent to it, the label
may be replaced by any index whatever provided a corresponding right-hand side is
supplied. \S\ref{sec:interp} replaces $d\in\{x,y,L,\dots\}$ by the discrete velocity
$q\in\{0,\dots,18\}$ and the right-hand side by $\bm X(\bm\delta_q)$. Nothing else in
the construction changes, which is why one factorisation continues to serve every case
(Remark~\ref{rem:shared}). The passivity of the superscript is precisely what makes the
extension available.
\end{remark}

\subsection{The same computation through the basis functions}

The example above bypassed \eqref{eq:wPsi}. Carrying the parameterisation through
exhibits its role. Adopt the Hermite construction in the probabilists' normalisation
of \S\ref{sec:abf}, write $a\equiv s/h_i$ and $\psi_a\equiv\psi(a)$, and use
$\mathrm{He}_1(z)=z$, $\mathrm{He}_2(z)=z^{2}-1$.

Here $p=2$ and $N_i=2$, so there are four numbers $W^{\alpha}_{ji}$, one for each pair
of a basis function and a neighbour. It is worth displaying them as an array, since the
roles of the two indices are the point of the exercise
(Remark~\ref{rem:indices}): \emph{both} basis functions are evaluated at \emph{both}
neighbours, and differ between them only through the argument.
\begin{center}
\begin{tabular}{@{}lcc@{}}
\toprule
& basis $\alpha{=}1$:\ $\mathrm{He}_1(z)=z$
& basis $\alpha{=}2$:\ $\mathrm{He}_2(z)=z^{2}-1$ \\
\midrule
neighbour $j_1$\ \ ($\tilde x=+a$) & $\psi_a\,a$ & $\psi_a\,(a^{2}-1)$ \\
neighbour $j_2$\ \ ($\tilde x=-a$) & $-\psi_a\,a$ & $\psi_a\,(a^{2}-1)$ \\
\bottomrule
\end{tabular}
\end{center}
\noindent The rows of this array are the vectors of \eqref{eq:wPsi},
\begin{equation}
\bm W_{j_1 i}=\psi_a\bigl[a,\ a^{2}-1\bigr]^{\mathsf T},
\qquad
\bm W_{j_2 i}=\psi_a\bigl[-a,\ a^{2}-1\bigr]^{\mathsf T},
\label{eq:exW}
\end{equation}
and the array itself is the matrix $\bm W$ of \S\ref{sec:structure}, of shape
$N_i\times p$. The sign difference between the rows in the first column, and its absence
in the second, is the parity of $\mathrm{He}_1$ and $\mathrm{He}_2$; it is what produces
the diagonal structure below.

Assembling $\bm M_i=\sum_j\bm X_{ji}\otimes\bm W_{ji}$ entrywise --- each entry being a
sum over neighbours, its position fixed by (which monomial, which basis function) ---
\begin{equation}
M_{11}=\underbrace{s\cdot\psi_a a}_{\text{from }j_1}
      +\underbrace{(-s)\cdot(-\psi_a a)}_{\text{from }j_2}=2sa\,\psi_a,
\qquad
M_{12}=\underbrace{s\cdot\psi_a(a^{2}{-}1)}_{\text{from }j_1}
      +\underbrace{(-s)\cdot\psi_a(a^{2}{-}1)}_{\text{from }j_2}=0,
\label{eq:exMentry}
\end{equation}
and likewise $M_{21}=0$ and $M_{22}=s^{2}\psi_a(a^{2}-1)$, so that
\begin{equation}
\bm M_i=\bm X^{\mathsf T}\bm W=
\begin{pmatrix}
2sa\,\psi_a & 0\\[0.2em]
0 & s^{2}\psi_a\bigl(a^{2}-1\bigr)
\end{pmatrix}.
\label{eq:exM}
\end{equation}
The off-diagonal entries vanish identically, in agreement with \eqref{eq:parity}: the
stencil is symmetric, and an entry pairing a monomial of one parity with a Hermite
function of the other integrates to zero. Solving \eqref{eq:moments} for
$\bm C=[1,0]^{\mathsf T}$ gives
$\bm\Psi=\bigl[(2sa\psi_a)^{-1},\,0\bigr]^{\mathsf T}$ and, through \eqref{eq:wPsi},
\begin{equation}
w_{j_1}=\psi_a a\cdot\frac{1}{2sa\psi_a}=\frac{1}{2s},
\qquad
w_{j_2}=-\psi_a a\cdot\frac{1}{2sa\psi_a}=-\frac{1}{2s},
\label{eq:excancel}
\end{equation}
identical to \eqref{eq:excd1}; the same cancellation occurs for
$\bm C=[0,1]^{\mathsf T}$. Neither $\psi$ nor $a$ survives. This is
Proposition~\ref{prop:degenerate} in the smallest possible instance: with $N_i=p$ the
parameterisation is a bijection, the selection it performs is vacuous, and the choice
of basis is immaterial.

\subsection{A degeneracy of the parameterisation that the problem does not possess}

Equation \eqref{eq:exM} is singular when $a=1$, that is when $h_i=s$, since the
second diagonal entry then vanishes. At that ratio the second-order operator cannot be
constructed through the ansatz, although the underlying moment conditions
\eqref{eq:excond} remain uniquely solvable, $\bm X$ being invertible by \eqref{eq:exX}.
The failure therefore belongs to the parameterisation and not to the problem.

This is the mechanism behind the requirement of \S\ref{sec:sampling} in its simplest
form. The statement of \cite[\S 3.2]{KLN2020} that the node distribution must
``adequately sample the ABFs'' is, in this example, the statement that $h_i/s$ must be
kept away from a value at which the chosen basis becomes blind to a moment it is
required to control. It is consistent that the operating ratios recommended
in~\cite[\S 5]{KL2021} for $m=2$ lie above unity; the example does not derive that
value, but it does show why some such condition is unavoidable, and why the condition
attaches to the basis rather than to the accuracy requirement.

\subsection{The underdetermined case}
\label{app:under}

Retain $m=2$ and add a third neighbour at $x=2s$, so that $N_i=3>p=2$. The moment
conditions for $\bm C=[1,0]^{\mathsf T}$ read
\begin{equation}
s\,(w_1-w_2+2w_3)=1,
\qquad
\tfrac{s^{2}}{2}\,(w_1+w_2+4w_3)=0,
\label{eq:exunder}
\end{equation}
whose general solution is the one-parameter family
\begin{equation}
w_3=t,
\qquad
w_1=\frac{1}{2s}-3t,
\qquad
w_2=-\frac{1}{2s}-t,
\qquad t\in\mathbb R .
\label{eq:exfamily}
\end{equation}
Every member reproduces polynomials of degree at most two exactly: on $\phi=x$ the
operator returns $s(w_1-w_2+2w_3)=1$, and on $\phi=x^{2}$ it returns
$s^{2}(w_1+w_2+4w_3)=0$, independently of $t$. The choice $t=0$ recovers
\eqref{eq:excd1}; $t=1/6s$ gives the equally consistent formula
$\bigl(0,\,-2/3s,\,1/6s\bigr)$.

The members are not, however, indistinguishable. Applying the operator to $\phi=x^{3}$,
the lowest degree not controlled by the moment conditions, gives
\begin{equation}
\mathcal L^{x}_i\bigl(x^{3}\bigr)=s^{3}\bigl(w_1-w_2+8w_3\bigr)=s^{2}+6s^{3}t ,
\label{eq:exleading}
\end{equation}
against an exact value of zero. The error is $O(s^{2})=O(s^{m})$ for every $t$, as the
order statement requires, but its coefficient depends linearly on $t$: the family shares
one order of accuracy and not one error constant. The precise
statement is therefore that the surplus degrees of freedom cannot change the
\emph{order}, only the coefficient of the leading error term, together with the
conditioning and stability properties of the resulting operator. It is this coefficient,
and not the order, that the resolving-power optimisations of~\cite{Broadley2025,%
Broadley2026} address.

Accuracy of a given order therefore does not distinguish the members of
\eqref{eq:exfamily}, and some further principle must. In LABFM that principle is the restriction of $\bm w$ to
$\operatorname{span}(\bm W)$, which selects one value of $t$; in the control formulation
of \S\ref{sec:mls} it is a weighted least-squares minimisation, which selects another.
The surplus degree of freedom is spent on conditioning and stability, as recorded in
Remark~\ref{rem:surplus}, and the requirement $N_i>p$ is not optional: the critical
neighbour counts of \cite[\S 3.2]{KLN2020} exceed the corresponding polynomial counts
at every order.

\subsection{The remaining assertions of \texorpdfstring{\cite[\S 3.1]{KLN2020}}{Sec. 3.1}}

Four statements close the section, and each may now be read directly.

\begin{enumerate}[label=(\roman*)]
\item \emph{Only the first $p$ elements of $\bm X$ and $\bm W$ are retained.} The
expansion \eqref{eq:momexpand} has infinitely many terms; finitely many can be
annihilated. The truncation point defines the method's order and fixes the size of
$\bm M_i$ through \eqref{eq:p}.
\item \emph{The leading error scales as $h^{\,m+1-l}$.} The first moment \emph{not}
prescribed is that of total degree $m+1$, and it multiplies a derivative of that order;
dividing by the $h^{-l}$ scaling of the weights for an $l$-th derivative gives the
quoted exponent. For $l=0$, the case of interest here, the exponent is $m+1$.
\item \emph{$\bm M_i$ and $\bm w^{d}$ depend on the anisotropy of the distribution,
yet the accuracy is maintained.} The diagonal form \eqref{eq:exM} is an artefact of
symmetry; on a disordered cloud $\bm M_i$ is full, and its off-diagonal entries measure
the extent to which each basis function contributes to moments other than its own. That
contamination is inverted by solving \eqref{eq:moments}, which is why the accuracy is
unaffected while the conditioning is not.
\item \emph{$\bm M_i$ is the same for every operator.} The matrix is assembled from the
stencil geometry alone; the operator enters only through the right-hand side. The
present work rests on this observation, and \S\ref{sec:interp} exploits it by supplying
a right-hand side that no longer corresponds to a derivative.
\end{enumerate}

\subsection{Relation to the operator of \texorpdfstring{\S\ref{sec:interp}}{Section 4}}

In the language of Remark~\ref{rem:distribution}, \cite[\S 3.1]{KLN2020} specifies the
moments of a derivative of a Dirac distribution centred at $\bm r_i$, and the entries of
$\bm C^{d}$ in \cite[Eq.~(5)]{KLN2020} are exactly those moments. The operator of
Definition~\ref{def:interp} changes nothing in the construction except the target: the
moments of a Dirac distribution translated to $\bm r_i+\bm\delta_q$, which are the
monomials evaluated at the displacement. Where \cite[\S 3.1]{KLN2020} sets almost every
moment to zero, the present right-hand side prescribes all of them to non-zero values;
the word ``elimination'' is then no longer descriptive, but the counting argument that
determines the order --- the first unprescribed moment is the leading error --- is
unchanged, and is what Theorem~\ref{thm:exact} formalises.

%========================================================================
\begin{thebibliography}{99}

\bibitem{KLN2020}
J.~R.~C. King, S.~J. Lind, A. Nasar,
\emph{High order difference schemes using the local anisotropic basis function method},
J. Comput. Phys. \textbf{415} (2020) 109549.
\url{https://doi.org/10.1016/j.jcp.2020.109549}.
\newline\emph{Cited throughout from the published version; local copy at}
\texttt{docs/literature/1.\,High order difference schemes using the local anisotropic
basis function method.pdf}\emph{. Section and equation numbers follow that version.}

\bibitem{KL2021}
J.~R.~C. King, S.~J. Lind,
\emph{High-order simulations of isothermal flows using the local anisotropic basis
function method (LABFM)},
J. Comput. Phys. \textbf{449} (2022) 110760.
\url{https://doi.org/10.1016/j.jcp.2021.110760}.
\newline\emph{Cited throughout from the published version; local copy at}
\texttt{docs/literature/2.\,High-order simulations of isothermal flows using the local
anisotropic basis function method (LABFM) pdf.pdf}\emph{. Section and equation numbers
follow that version.}

\bibitem{Broadley2025}
H. Broadley, S.~J. Lind, J.~R.~C. King,
\emph{Improving the accuracy of meshless methods via resolving power optimisation using
multiple kernels}, arXiv:2510.20365 (2025).
\url{https://arxiv.org/abs/2510.20365}.

\bibitem{Broadley2026}
H.~M. Broadley, J.~R.~C. King, S.~J. Lind,
\emph{Compact LABFM: a framework for meshless methods with spectral-like resolving
power}, arXiv:2603.11668 (2026).
\url{https://arxiv.org/abs/2603.11668}.

\bibitem{Feng2025}
R. Feng, J.~R.~C. King, S.~J. Lind,
\emph{A p-adaptive high-order mesh-free framework for fluid simulations in complex
geometries}, arXiv:2511.21224 (2025).
\url{https://arxiv.org/abs/2511.21224}.

\bibitem{Lancaster1981}
P. Lancaster, K. Salkauskas,
\emph{Surfaces generated by moving least squares methods},
Math. Comp. \textbf{37} (1981) 141--158.
\url{https://doi.org/10.1090/S0025-5718-1981-0616367-1}.

\bibitem{Mirzaei2012}
D. Mirzaei, R. Schaback, M. Dehghan,
\emph{On generalized moving least squares and diffuse derivatives},
IMA J. Numer. Anal. \textbf{32} (2012) 983--1000.
\url{https://doi.org/10.1093/imanum/drr030}.

\bibitem{Trask2016}
N. Trask, M. Maxey, X. Hu,
\emph{Compact moving least squares: an optimization framework for generating
high-order compact meshless discretizations},
J. Comput. Phys. \textbf{326} (2016) 596--611.
\url{https://doi.org/10.1016/j.jcp.2016.08.045}.

\bibitem{FornbergFlyer2015}
B. Fornberg, N. Flyer,
\emph{Fast generation of 2D node distributions for mesh-free PDE discretizations},
Comput. Math. Appl. \textbf{69} (2015) 531--544.

\bibitem{Musavi2015}
S.~H. Musavi, M. Ashrafizaadeh,
\emph{Meshless lattice Boltzmann method for the simulation of fluid flows},
Phys. Rev. E \textbf{91} (2015) 023310.
\url{https://doi.org/10.1103/PhysRevE.91.023310}.

\bibitem{LinWuZhang2019}
X. Lin, J. Wu, T. Zhang,
\emph{A mesh-free radial basis function-based semi-Lagrangian lattice Boltzmann method
for incompressible flows},
Int. J. Numer. Methods Fluids (2019).
\url{https://doi.org/10.1002/fld.4749}.

\bibitem{Strzelczyk2023}
D. Strzelczyk, M. Matyka,
\emph{Study of the convergence of the meshless lattice Boltzmann method in
Taylor--Green and annular channel flows}, arXiv:2306.01525 (2023).
\url{https://arxiv.org/abs/2306.01525}.

\bibitem{Strzelczyk2024}
D. Strzelczyk, M. Matyka,
\emph{Modeling inertial flows with meshless lattice Boltzmann method},
arXiv:2312.14789 (2024).
\url{https://arxiv.org/abs/2312.14789}.

\bibitem{WangYang2025}
X. Wang, S. Yang,
\emph{A semi-Lagrangian meshfree lattice Boltzmann method for incompressible two-phase
flows}, J. Comput. Phys. (2025).
\url{https://doi.org/10.1016/j.jcp.2025.113727}.

\bibitem{Maidenberg2026}
A. Maidenberg, Y. Chen,
\emph{A meshless discrete Boltzmann solver with the radial basis function finite
difference scheme for fluid flows with Dirichlet boundary conditions},
J. Fluids Eng. \textbf{148} (2026) 021501.
\url{https://doi.org/10.1115/1.4069638}.

\bibitem{Kramer2017}
A. Kr\"amer, K. K\"ullmer, D. Reith, W. Joppich, H. Foysi,
\emph{Semi-Lagrangian off-lattice Boltzmann method for weakly compressible flows},
Phys. Rev. E \textbf{95} (2017) 023305.
\url{https://doi.org/10.1103/PhysRevE.95.023305}.

\bibitem{Wilde2020}
D. Wilde, A. Kr\"amer, D. Reith, H. Foysi,
\emph{Semi-Lagrangian lattice Boltzmann method for compressible flows},
Phys. Rev. E \textbf{101} (2020) 053306.
\url{https://doi.org/10.1103/PhysRevE.101.053306}.

\bibitem{Wilde2021}
D. Wilde, A. Kr\"amer, D. Reith, H. Foysi,
\emph{High-order semi-Lagrangian kinetic scheme for compressible turbulence},
Phys. Rev. E \textbf{104} (2021) 025301.
\url{https://doi.org/10.1103/PhysRevE.104.025301};
preprint \url{https://arxiv.org/abs/2012.05537}.

\bibitem{Spelten2024}
P. Spelten, D. Wilde, M.~C. Bedrunka, D. Reith, H. Foysi,
\emph{Supersonic shear and wall-bounded flows with body-fitted meshes using the
semi-Lagrangian lattice Boltzmann method: boundary schemes and applications},
arXiv:2412.09051 (2024).
\url{https://arxiv.org/abs/2412.09051}.

\bibitem{Jin2025}
B.-X. Jin, S.-W. Feng, Y.-H. Chiu, C.-A. Lin,
\emph{Direct numerical simulations of turbulent channel and duct flows using
interpolation-based lattice Boltzmann method},
Flow Turbul. Combust. (2025).
\url{https://doi.org/10.1007/s10494-025-00689-w}.

\bibitem{Breuer2009}
M. Breuer, N. Peller, C. Rapp, M. Manhart,
\emph{Flow over periodic hills --- numerical and experimental study in a wide range of
Reynolds numbers}, Comput. Fluids \textbf{38} (2009) 433--457.

\end{thebibliography}

\end{document}
